% This chapter has been modified on 6/4/05.
%\setcounter{chapter}{3}

\chapter{Peluang Bersyarat}\label{chp 4}  

\section{Peluang Bersyarat Diskret}\label{sec 4.1} 

\subsection*{Peluang Bersyarat}
\par
Dalam bagian ini kita mengajukan dan menjawab pertanyaan berikut.  Misalkan kita menetapkan
suatu fungsi distribusi pada sebuah ruang sampel, lalu mengetahui bahwa kejadian $E$ telah
terjadi.  Bagaimana seharusnya kita mengubah peluang kejadian-kejadian yang tersisa?  Peluang
baru untuk suatu kejadian $F$ akan kita sebut \emx {peluang
bersyarat\index{peluang!bersyarat}\index{peluang bersyarat} $F$ dengan syarat $E$} dan
kita lambangkan dengan
$P(F|E)$.

\begin{example}\label{exam 4.1}
Suatu percobaan terdiri atas pelemparan sebuah dadu satu kali.  Misalkan $X$ adalah hasilnya.  Misalkan $F$ adalah kejadian
$\{X = 6\}$, dan misalkan $E$ adalah kejadian $\{X > 4\}$.  Kita menetapkan
fungsi distribusi
$m(\omega) = 1/6$ untuk $\omega = 1, 2, \ldots , 6$.  Jadi, $P(F) =
1/6$.  Sekarang misalkan dadu dilempar dan kita diberi tahu bahwa kejadian $E$ telah terjadi.  
Dengan demikian hanya tersisa dua hasil yang mungkin: 5 dan~6.  Tanpa informasi lain, 
kita tetap menganggap kedua hasil ini berpeluang sama, sehingga 
peluang $F$ menjadi 1/2, jadi $P(F|E) = 1/2$.
\end{example}

\begin{example}\label{exam 4.1.5}
Dalam Tabel Hayat (lihat Lampiran~C), kita mendapati bahwa dalam populasi 100{,}000 perempuan, 89.835\%
diperkirakan hidup hingga usia 60 tahun, sedangkan 57.062\% diperkirakan hidup hingga usia 80 tahun.  Dengan syarat seorang perempuan telah berusia
60 tahun, berapakah peluang bahwa ia hidup hingga usia 80 tahun?
\par
Ini adalah contoh peluang bersyarat.  Dalam kasus ini, ruang sampel semula
dapat dipandang sebagai himpunan 100{,}000 perempuan.  Kejadian $E$ dan $F$ masing-masing adalah subhimpunan
ruang sampel yang terdiri atas semua perempuan yang hidup sekurang-kurangnya 60 tahun dan sekurang-kurangnya 80 tahun.
Kita menganggap $E$ sebagai ruang sampel baru dan mencatat bahwa $F$ merupakan subhimpunan $E$. 
Jadi, ukuran $E$ adalah 89{,}835 dan ukuran $F$ adalah 57{,}062.  Maka, peluang yang
ditanyakan sama dengan $57{,}062/89{,}835 = .6352$.  Dengan demikian, seorang perempuan berusia 60 tahun memiliki peluang 63.52\% untuk
hidup hingga usia 80 tahun.
\end{example}


\begin{example}\label{exam 4.2}
Perhatikan kembali contoh pemungutan suara dari Bagian~\ref{sec 1.2}: tiga kandidat A, B,
dan C mencalonkan diri.  Kita memutuskan bahwa A dan B memiliki peluang yang sama untuk
menang, sedangkan peluang C untuk menang hanya 1/2 peluang A.  Misalkan $A$ adalah kejadian ``A
menang," $B$ adalah kejadian ``B menang," dan $C$ adalah kejadian ``C menang."  Karena itu, kita menetapkan
peluang $P(A) = 2/5$, $P(B) = 2/5$, dan $P(C) = 1/5$.  
\par
Misalkan sebelum pemilihan berlangsung, $A$ mengundurkan diri.  Seperti dalam Contoh~\ref{exam
4.1}, wajar jika kita menetapkan peluang baru untuk kejadian $B$ dan $C$ yang
sebanding dengan peluang semula.  Dengan demikian, kita memperoleh $P(B|~A) = 2/3$, dan
$P(C|~A) = 1/3$.  Penting untuk dicatat bahwa setiap kali kita menetapkan peluang pada kejadian di dunia nyata,
distribusi yang dihasilkan hanya berguna jika kita memperhitungkan semua informasi yang
relevan.  Dalam contoh ini, kita mungkin mengetahui bahwa sebagian besar pemilih yang mendukung $A$ akan memilih
$C$ jika $A$ tidak lagi ikut serta.  Hal ini jelas membuat peluang
$C$ menang lebih besar daripada nilai 1/3 yang ditetapkan di atas.
\end{example}

Dalam contoh-contoh ini kita menetapkan suatu fungsi distribusi, lalu memperoleh informasi
baru yang menentukan ruang sampel baru, yang terdiri atas hasil-hasil yang
masih mungkin, dan membuat kita menetapkan fungsi distribusi baru pada
ruang tersebut.

Kita hendak memformalkan prosedur yang dilakukan dalam contoh-contoh tersebut.  Misalkan
$\Omega = \{\omega_1,\omega_2,\dots,\omega_r\}$ adalah ruang sampel semula
yang telah diberi fungsi distribusi $m(\omega_j)$.  Misalkan kita mengetahui bahwa
kejadian $E$ telah terjadi.  Kita hendak menetapkan fungsi distribusi baru
$m(\omega_j|E)$ pada $\Omega$ untuk mencerminkan fakta ini.  Jelas bahwa jika suatu titik sampel
$\omega_j$ tidak berada dalam $E$, kita menginginkan $m(\omega_j|E) = 0$.  Selain itu, jika tidak ada
informasi yang menyatakan sebaliknya, wajar untuk menganggap bahwa
peluang bagi $\omega_k$ dalam $E$ harus mempertahankan perbandingan relatif yang sama
seperti sebelum kita mengetahui bahwa $E$ telah terjadi.  Untuk itu kita mensyaratkan
bahwa
$$
m(\omega_k|E) = cm(\omega_k)
$$
untuk semua $\omega_k$ dalam $E$, dengan $c$ suatu konstanta positif.  Namun, kita juga harus
memiliki
$$
\sum_E m(\omega_k|E) = c\sum_E m(\omega_k) = 1\ .
$$
Jadi,
$$
c = \frac 1{\sum_E m(\omega_k)} = \frac 1{P(E)}\ .
$$
(Perhatikan bahwa hal ini mengharuskan kita mengasumsikan $P(E) > 0$.)  Dengan demikian, kita
mendefinisikan
$$
m(\omega_k|E) = \frac {m(\omega_k)}{P(E)}
$$
untuk $\omega_k$ dalam $E$.  Distribusi baru ini akan kita sebut \emx {distribusi
bersyarat}\index{distribusi bersyarat} dengan syarat $E$.  Untuk suatu kejadian umum $F$, diperoleh
$$
P(F|E) = \sum_{F \cap E} m(\omega_k|E) = \sum_{F \cap E}\frac {m(\omega_k)}{P(E)} = 
\frac {P(F \cap E)}{P(E)}\ .
$$

Kita menyebut $P(F|E)$ sebagai \emx {peluang bersyarat terjadinya $F$ dengan syarat $E$ terjadi,}
dan menghitungnya dengan rumus
$$
P(F|E) = \frac {P(F \cap E)}{P(E)}\ .
$$

\begin{example}(Lanjutan Contoh~\ref{exam 4.1}){\label{exam 4.3}}
Mari kembali ke contoh pelemparan dadu.  Ingat bahwa $F$ adalah kejadian $X = 6$, dan
$E$ adalah kejadian $X > 4$.  Perhatikan bahwa $E \cap F$ adalah kejadian $F$.  Jadi, rumus di atas memberikan
\begin{eqnarray*}
P(F|E) & = & \frac {P(F \cap E)}{P(E)} \\
& = & \frac {1/6}{1/3} \\
& = & \frac 12\ ,\\
\end{eqnarray*}
sesuai dengan perhitungan sebelumnya.
\end{example}

\begin{example}\label{exam 4.4}
Kita memiliki dua guci, I dan II.  Guci I berisi 2 bola hitam dan 3 bola putih. 
Guci II berisi 1 bola hitam dan 1 bola putih.  Sebuah guci dipilih secara acak dan sebuah bola
dipilih secara acak dari guci tersebut.  Kita dapat menyajikan ruang sampel percobaan ini sebagai lintasan-lintasan
pada sebuah pohon seperti ditunjukkan dalam Gambar~\ref{fig 4.1}. Peluang yang ditetapkan pada lintasan-lintasan tersebut
juga ditampilkan.  
\par
Misalkan $B$ adalah kejadian ``sebuah bola hitam terambil," dan $I$ adalah kejadian ``guci I
terpilih."  Bobot cabang 2/5, yang ditunjukkan pada salah satu cabang dalam gambar, kini dapat
ditafsirkan sebagai peluang bersyarat $P(B|I)$.

\putfig{4.5truein}{PSfig4-1}{Diagram pohon.}{fig 4.1}

\par
Misalkan kita hendak menghitung $P(I|B)$.  Dengan menggunakan rumus tersebut, kita memperoleh
\begin{eqnarray*}
P(I|B) & = & \frac {P(I \cap B)}{P(B)} \\
       & = & \frac {P(I \cap B)}{P(B \cap I) + P(B \cap II)} \\
       & = & \frac {1/5}{1/5 + 1/4} = \frac 49\ .
\end{eqnarray*}
\end{example}

\subsection*{Peluang Bayes}

Ukuran pohon semula memberikan peluang terambilnya bola dengan
warna tertentu, dengan syarat guci yang dipilih diketahui.  Kita baru saja menghitung \emx {peluang
balikan} bahwa guci tertentu telah dipilih, dengan syarat warna bolanya diketahui. 
Peluang balikan semacam ini disebut \emx {peluang Bayes}\index{peluang
Bayes}\index{peluang!Bayes} dan dapat diperoleh dengan rumus yang akan kita kembangkan
kemudian.  Peluang Bayes juga dapat diperoleh hanya dengan menyusun ukuran pohon untuk
percobaan dua tahap yang dilakukan dalam urutan terbalik.  Pohon ini ditampilkan dalam Gambar~\ref{fig 4.2}.

\putfig{4.5truein}{PSfig4-2}{Diagram pohon terbalik.}{fig 4.2}

Lintasan-lintasan pada pohon terbalik berkorespondensi satu-satu dengan lintasan pada pohon maju,
karena masing-masing bersesuaian dengan hasil tertentu dari percobaan, sehingga diberi
peluang yang sama.  Dari pohon maju, kita memperoleh bahwa peluang terambilnya bola hitam adalah
$$
\frac 12 \cdot \frac 25 + \frac 12 \cdot \frac 12 = \frac 9{20}\ .
$$

Peluang untuk cabang-cabang pada tingkat kedua diperoleh dengan pembagian
sederhana.  Sebagai contoh, jika $x$ adalah peluang yang akan ditetapkan pada
cabang teratas di tingkat kedua, kita harus mempunyai
$$
\frac 9{20} \cdot x = \frac 15
$$
atau $x = 4/9$.  Jadi, $P(I|B) = 4/9$, sesuai dengan
perhitungan kita sebelumnya.  Pohon terbalik itu kemudian menampilkan semua peluang
balikan, atau peluang Bayes.  

\begin{example}\label{exam 4.5}
Sekarang kita membahas suatu masalah yang disebut masalah \emx {Monty Hall}\index{masalah Monty Hall}.   Masalah ini
telah lama menjadi masalah favorit, tetapi kembali populer setelah Craig Whitaker\index{WHITAKER,
C.} mengirim surat kepada Marilyn vos Savant\index{vos SAVANT, M.} untuk dibahas dalam kolomnya di \emx {Parade
Magazine}.\footnote{Marilyn vos Savant, Ask Marilyn, \emx {Parade Magazine}, 9 
September;  2 December;  17  February 1990, dicetak ulang dalam Marilyn vos Savant, \emx {Ask
Marilyn}, St. Martins, New York, 1992.}  Persoalannya dapat dirumuskan secara mandiri sebagai berikut.
Terdapat tiga pintu: satu menyembunyikan mobil\index{mobil} dan dua lainnya kambing\index{kambing}.
Peserta memilih pintu~1. Pembawa acara, yang mengetahui letak hadiah, lalu membuka pintu~3 dan
memperlihatkan seekor kambing, kemudian menawarkan pintu~2 sebagai pengganti. Apakah beralih pintu
menguntungkan? Uraian ini tidak menerjemahkan atau mereproduksi ungkapan surat pembaca.
\par
Marilyn memberikan solusi yang menyimpulkan bahwa Anda sebaiknya beralih, dan jika melakukannya, peluang Anda untuk
menang adalah 2/3.  Beberapa pembaca yang marah, sebagian di antaranya menyatakan memiliki gelar PhD dalam
matematika, mengatakan bahwa hal itu tidak masuk akal karena setelah Monty menyingkirkan satu pintu, hanya ada dua
pintu yang mungkin dan masing-masing seharusnya tetap memiliki peluang 1/2, sehingga beralih
tidak memberikan keuntungan.  Marilyn mempertahankan solusinya dan mendorong para pembacanya untuk menyimulasikan permainan tersebut serta menarik
kesimpulan sendiri.  Kami juga mendorong pembaca untuk melakukannya (lihat
Latihan~\ref{exer 4.1.9.6}).
\par
Pembaca lain mengeluhkan bahwa Marilyn belum menggambarkan masalah itu secara lengkap.  Secara khusus,
cara pengambilan keputusan tertentu selama permainan tidak dijelaskan.  Aspek
masalah ini akan dibahas dalam Bagian~\ref{sec 4.3}.  Kita akan mengasumsikan bahwa mobil ditempatkan
di balik sebuah pintu dengan melempar dadu bersisi tiga, sehingga ketiga pilihan berpeluang sama.  Monty
mengetahui letak mobil dan selalu membuka pintu yang di baliknya ada kambing.  Terakhir, kita mengasumsikan bahwa jika 
Monty dapat memilih di antara dua pintu (yaitu, peserta telah memilih pintu yang di baliknya ada mobil), ia
memilih masing-masing pintu dengan peluang 1/2.  Marilyn
jelas mengharapkan para pembacanya mengasumsikan bahwa permainan berlangsung dengan cara ini.
\par
Seperti kebanyakan paradoks semu, masalah ini dapat diselesaikan melalui analisis yang cermat.
Kita mulai dengan menjelaskan pertanyaan terkait yang lebih sederhana.  Kita katakan bahwa seorang peserta menggunakan
strategi ``bertahan" jika ia memilih sebuah pintu dan, ketika ditawari kesempatan untuk beralih ke pintu lain,
menolak (yaitu, ia bertahan pada pilihan semula).  Demikian pula, kita katakan bahwa
peserta menggunakan strategi ``beralih" jika ia memilih sebuah pintu dan, ketika ditawari kesempatan untuk beralih ke
pintu lain, menerima tawaran itu.  Sekarang misalkan seorang peserta memutuskan sejak awal untuk memakai
strategi ``bertahan."  Satu-satunya tindakannya dalam kasus ini adalah memilih sebuah pintu (dan menolak tawaran untuk
beralih, jika tawaran diberikan).  Berapakah peluang ia memenangkan sebuah mobil?  Pertanyaan yang sama
dapat diajukan untuk strategi ``beralih."  
\par
Dengan strategi ``bertahan," seorang peserta akan memenangkan mobil dengan peluang 1/3, karena dalam 1/3
kasus mobil berada di balik pintu yang dipilihnya.  Di pihak lain, jika seorang peserta memakai
strategi ``beralih," ia akan menang setiap kali mobil tidak berada di balik pintu yang semula
dipilihnya, yang terjadi dalam 2/3 kasus.  
\par
Analisis yang sangat sederhana ini, meskipun benar, belum sepenuhnya menyelesaikan masalah yang diajukan Craig.  Craig
menanyakan peluang bersyarat bahwa Anda menang jika beralih, dengan syarat Anda telah memilih pintu 1
dan Monty telah memilih pintu 3.  Untuk menyelesaikan masalah ini, kita menyusun masalah sebelum
memperoleh informasi tersebut, lalu menghitung peluang bersyarat berdasarkan informasi itu. 
Proses ini berlangsung dalam beberapa tahap; mobil ditempatkan di balik sebuah pintu, peserta
memilih sebuah pintu, dan akhirnya Monty membuka sebuah pintu.  Karena itu, wajar untuk menganalisisnya dengan ukuran
pohon.  Di sini kita membuat asumsi tambahan bahwa jika Monty dapat memilih di antara dua pintu (yaitu,
peserta telah memilih pintu yang di baliknya ada mobil), ia memilih masing-masing pintu dengan peluang
1/2.  Asumsi-asumsi yang kita buat menentukan peluang cabang, yang kemudian menentukan
ukuran pohon. Pohon dan ukuran pohon yang dihasilkan ditampilkan dalam Gambar~\ref{fig 4.4}.  Mungkin terasa menggoda untuk
memperkecil pohon dengan membuat asumsi tertentu, misalnya: ``Tanpa mengurangi keumuman, kita
akan mengasumsikan bahwa peserta selalu memilih pintu 1."  Demi kejelasan, kita memilih untuk tidak membuat
asumsi semacam itu.  

\putfig{4.5truein}{PSfig4-0}{Masalah Monty Hall.}{fig 4.4}
\par
Sekarang informasi yang diberikan, yaitu peserta memilih pintu 1 dan Monty memilih pintu 3, berarti
hanya dua lintasan pada pohon yang mungkin (lihat Gambar~\ref{fig 4.4.5}).  
\putfig{4.5truein}{PSfig4-4-5}{Peluang bersyarat untuk masalah Monty Hall.}{fig 4.4.5}
Pada salah satu lintasan, mobil berada di balik pintu 1, sedangkan pada lintasan lainnya mobil berada di balik pintu 2.  
Lintasan dengan mobil di balik pintu 2 berpeluang dua kali lebih besar daripada lintasan dengan mobil di balik pintu 1. 
Jadi, peluang bersyarat bahwa mobil berada di balik pintu 2 adalah 2/3 dan bahwa mobil berada 
di balik pintu 1 adalah 1/3. Karena itu, jika Anda beralih, peluang Anda memenangkan mobil adalah 2/3, seperti yang dikatakan Marilyn.
\par
Pada titik ini, pembaca mungkin mengira bahwa kedua masalah di atas sama karena jawabannya
sama.  Ingat bahwa dalam masalah semula kita mengasumsikan bahwa jika peserta
memilih pintu dengan mobil di baliknya, sehingga Monty dapat memilih di antara dua pintu, Monty memilih masing-masing pintu dengan
peluang 1/2.  Sekarang misalkan sebagai gantinya, ketika Monty dapat memilih, ia memilih pintu
dengan nomor lebih besar dengan peluang 3/4.  Dalam masalah ``beralih" versus ``bertahan,"
peluang menang dengan strategi ``beralih" tetap 2/3.  Namun, dalam
masalah semula, jika peserta beralih, ia menang dengan peluang 4/7.  Pembaca dapat
memeriksanya dengan mencatat bahwa hanya dua lintasan yang sama seperti sebelumnya yang mungkin pada pohon. 
Lintasan yang menghasilkan kemenangan jika peserta beralih memiliki peluang 1/3, sedangkan lintasan yang
menghasilkan kekalahan jika peserta beralih memiliki peluang 1/4.
\end{example}


\subsection*{Kejadian Bebas}     

Sering kali, pengetahuan bahwa suatu kejadian $E$ telah terjadi tidak memengaruhi
peluang terjadinya kejadian lain $F$, yaitu $P(F|E) = P(F)$.  Dalam
kasus ini, kita juga mengharapkan persamaan
$P(E|F) = P(E)$ berlaku.  Bahkan (lihat Latihan~\ref{exer 4.1.1}), 
masing-masing persamaan mengakibatkan persamaan yang lain.  Jika persamaan-persamaan ini berlaku, kita dapat mengatakan
bahwa $F$ \emx {bebas} dari $E$.  Sebagai contoh, kita tidak mengharapkan
pengetahuan tentang hasil lemparan pertama sebuah koin mengubah peluang
yang kita tetapkan pada hasil-hasil yang mungkin dari lemparan kedua; dengan kata lain, kita
tidak mengharapkan lemparan kedua bergantung pada lemparan pertama.  Gagasan ini
diformalkan dalam definisi kejadian bebas berikut.

\begin{definition}\label{def 4.2}
Misalkan $E$ dan $F$ adalah dua kejadian.  Kita mengatakan bahwa keduanya \emx {bebas}\index{kejadian!bebas}\index{kebebasan
kejadian} jika 1) kedua kejadian memiliki peluang positif dan
$$
P(E|F) = P(E)\ {\rm dan}\ P(F|E) = P(F)\ ,
$$
atau 2) sekurang-kurangnya salah satu kejadian memiliki peluang 0.
\end{definition}

Seperti dicatat di atas, jika $P(E)$ dan $P(F)$ keduanya positif, masing-masing persamaan di atas
mengakibatkan persamaan yang lain. Karena itu, untuk menentukan apakah dua kejadian bebas, hanya satu dari
persamaan tersebut yang perlu diperiksa (lihat Latihan~\ref{exer 4.1.1}).
\par
Teorema berikut memberikan cara lain untuk memeriksa kebebasan.

\begin{theorem}\label{thm 4.1}
Dua kejadian $E$ dan $F$ bebas jika dan hanya jika 
$$P(E\cap F) = P(E)P(F)\ .$$
\proof
Jika salah satu kejadian memiliki peluang 0, kedua kejadian itu bebas dan persamaan di atas 
berlaku, sehingga teorema benar dalam kasus ini.  Oleh karena itu, selanjutnya kita dapat mengasumsikan bahwa kedua kejadian memiliki 
peluang positif. Asumsikan bahwa $E$ dan $F$ bebas.  Maka $P(E|F) = P(E)$, sehingga
\begin{eqnarray*}
P(E\cap F) & = & P(E|F)P(F) \\
       & = & P(E)P(F)\ .
\end{eqnarray*}

Selanjutnya asumsikan bahwa $P(E\cap F) = P(E)P(F)$.  Maka
$$
P(E|F) = \frac {P(E \cap F)}{P(F)} = P(E)\ .
$$
Selain itu,
$$
P(F|E) = \frac {P(F \cap E)}{P(E)} = P(F)\ .
$$
Jadi, $E$ dan $F$ bebas.
\end{theorem}

\begin{example}\label{exam 4.6}
Misalkan kita memiliki sebuah koin yang menghasilkan kepala dengan peluang $p$ dan ekor dengan peluang
$q$.  Sekarang misalkan koin ini dilempar dua kali.  Dengan menggunakan penafsiran frekuensi atas peluang, wajar untuk menetapkan
pada hasil $(H,H)$ peluang $p^2$, pada hasil $(H, T)$ peluang $pq$, dan seterusnya.  
Misalkan $E$ adalah kejadian munculnya kepala pada lemparan pertama dan $F$ adalah kejadian munculnya ekor pada
lemparan kedua.  Sekarang kita akan memeriksa bahwa dengan penetapan peluang di atas, kedua kejadian ini
bebas, seperti yang diharapkan.  Kita mempunyai $P(E) = p^2 + pq = p$, $P(F) = pq + q^2 = q$.  Terakhir, $P(E\cap
F) = pq$, sehingga $P(E\cap F) = \linebreak[4] P(E)P(F)$.  
\end{example}

\begin{example}\label{exam 4.7}
Sering kali, tetapi tidak selalu, kita dapat melihat secara intuitif kapan dua kejadian
bebas.  Dalam Contoh~\ref{exam 4.6}, misalkan $A$ adalah kejadian ``lemparan
pertama menghasilkan kepala" dan $B$ adalah kejadian ``kedua hasil sama." Maka
$$
P(B|A) = \frac {P(B \cap A)}{P(A)} = \frac {P\{\mbox {HH}\}}{P\{\mbox {HH,HT}\}} =
\frac {1/4}{1/2} = \frac 12 = P(B).
$$
Jadi, $A$ dan $B$ bebas, meskipun hasil tersebut tidak begitu jelas secara intuitif.
\end{example}

\begin{example}\label{exam 4.8}
Terakhir, mari kita berikan contoh dua kejadian yang tidak bebas. 
Dalam Contoh~\ref{exam 4.6}, misalkan $I$ adalah kejadian ``kepala pada lemparan pertama" dan $J$
adalah kejadian ``muncul dua kepala."  Maka $P(I) = 1/2$ dan $P(J) = 1/4$.  Kejadian
$I \cap J$ adalah kejadian ``kepala pada kedua lemparan" dan memiliki peluang $1/4$. 
Jadi, $I$ dan $J$ tidak bebas karena $P(I)P(J) = 1/8 \ne P(I \cap J)$.
\end{example}

Kita dapat memperluas konsep kebebasan pada sebarang himpunan berhingga kejadian
$A_1$,~$A_2$, \dots,~$A_n$.

\begin{definition}\label{def 4.3}
Suatu himpunan kejadian $\{A_1,\ A_2,\ \ldots,\ A_n\}$ disebut \emx {saling
bebas}\index{kebebasan kejadian!saling}\index{kejadian saling bebas} jika untuk setiap
subhimpunan
$\{A_i,\ A_j,\ldots,\ A_m\}$ dari kejadian-kejadian tersebut berlaku
$$
P(A_i \cap A_j \cap\cdots\cap A_m) = P(A_i)P(A_j)\cdots P(A_m),
$$
atau secara ekuivalen, jika untuk sebarang barisan $\bar A_1$,~$\bar A_2$,
\dots,~$\bar A_n$ dengan $\bar A_j = A_j$ atau $\tilde A_j$,
$$
P(\bar A_1 \cap \bar A_2 \cap\cdots\cap \bar A_n) =
P(\bar A_1)P(\bar A_2)\cdots P(\bar A_n).
$$
(Untuk bukti keekuivalenan dalam kasus $n = 3$, 
lihat Latihan~\ref{exer 4.1.31}.)
\end{definition}
Dengan terminologi ini, setiap barisan  $(\mbox S,\mbox S,\mbox F,\mbox
F, \mbox S, \dots,\mbox S)$ dari hasil-hasil yang mungkin dalam suatu proses percobaan Bernoulli membentuk barisan kejadian yang saling
bebas.
\par
Wajar untuk bertanya: Jika setiap pasangan dalam suatu himpunan kejadian bebas, apakah keseluruhan himpunan itu
saling bebas?  Jawabannya adalah \emx {belum tentu,} dan sebuah contoh diberikan dalam
Latihan~\ref{exer 4.1.7}.
\par
Penting untuk dicatat bahwa pernyataan 
$$
P(A_1 \cap A_2 \cap \cdots \cap A_n) = P(A_1)P(A_2) \cdots P(A_n)
$$
tidak mengakibatkan bahwa kejadian-kejadian $A_1$,~$A_2$, \dots,~$A_n$ saling bebas (lihat
Latihan~\ref{exer 4.1.8}).

\subsection*{Fungsi Distribusi Gabungan dan Kebebasan Peubah Acak}

Sering kali, ketika suatu percobaan dilakukan, beberapa besaran berbeda yang
berkaitan dengan hasilnya diselidiki.  

\begin{example}\label{exam 4.91}
Misalkan kita melempar sebuah koin tiga kali.  Peubah acak dasar ${\bar X}$ yang berkaitan
dengan percobaan ini memiliki delapan hasil yang mungkin, yaitu tripel terurut yang terdiri
atas H dan T.  Kita juga dapat mendefinisikan peubah acak $X_i$, untuk $i = 1, 2, 3$, sebagai
hasil lemparan ke-$i$.  Jika koin itu adil, kita harus menetapkan peluang
1/8 pada masing-masing dari delapan hasil yang mungkin.  Jadi, fungsi distribusi $X_1$,
$X_2$, dan $X_3$ identik; dalam setiap kasus fungsi itu didefinisikan oleh $m(H) = m(T) = 1/2$.
\end{example}

Jika kita memiliki beberapa peubah acak $X_1, X_2, \ldots, X_n$ yang berkaitan dengan suatu
percobaan, kita dapat mempertimbangkan peubah acak gabungan ${\bar X} = (X_1, X_2,
\ldots, X_n)$ yang didefinisikan dengan mengambil suatu hasil $\omega$ dari percobaan dan menuliskan,
sebagai sebuah tupel-$n$, $n$ hasil yang bersesuaian untuk peubah acak $X_1, X_2, \ldots, X_n$. 
Jadi, jika himpunan hasil yang mungkin bagi peubah acak $X_i$ adalah himpunan $R_i$, maka
himpunan hasil yang mungkin bagi peubah acak gabungan ${\bar X}$ adalah hasil kali Kartesius
dari himpunan-himpunan $R_i$, yaitu himpunan semua tupel-$n$ hasil yang mungkin bagi peubah-peubah $X_i$.

\begin{example}(Lanjutan Contoh~\ref{exam 4.91}){\label{exam 4.92}}
Dalam contoh pelemparan koin di atas, misalkan $X_i$ menyatakan hasil lemparan ke-$i$.  Maka
peubah acak gabungan ${\bar X} = (X_1, X_2, X_3)$ memiliki delapan hasil yang mungkin.
\par
Sekarang misalkan kita mendefinisikan $Y_i$, untuk $i = 1, 2, 3$, sebagai banyaknya kepala yang muncul dalam
$i$ lemparan pertama.  Maka hasil yang mungkin bagi $Y_i$ adalah $\{0, 1, \ldots, i\}$, sehingga sekilas
himpunan hasil yang mungkin bagi peubah acak gabungan ${\bar Y} = (Y_1,
Y_2, Y_3)$ seharusnya merupakan himpunan
$$\{(a_1, a_2, a_3)\ :\ 0 \le a_1 \le 1, 0 \le a_2 \le 2, 0 \le a_3 \le 3\}\ .$$
Namun, hasil $(1, 0, 1)$ tidak mungkin terjadi karena harus berlaku $a_1 \le a_2 \le a_3$.  Penyelesaian
masalah ini adalah mendefinisikan peluang hasil $(1, 0, 1)$ sebagai 0.  Selain itu, harus
berlaku $a_{i+1} - a_i \le 1$ untuk $i = 1, 2$.
\par
Sekarang kita menggambarkan penetapan peluang pada berbagai hasil bagi peubah acak gabungan
${\bar X}$ dan ${\bar Y}$.  Dalam kasus pertama, masing-masing dari
delapan hasil harus diberi peluang 1/8 karena kita mengasumsikan koin yang digunakan
adil.  Dalam kasus kedua, karena $Y_i$ memiliki $i+1$ hasil yang mungkin, himpunan hasil yang
mungkin berukuran 24.    Hanya delapan dari 24 hasil tersebut yang benar-benar dapat terjadi, yaitu hasil yang
memenuhi $a_1 \le a_2 \le a_3$.  Masing-masing hasil ini bersesuaian dengan tepat satu
hasil peubah acak ${\bar X}$, sehingga wajar untuk menetapkan peluang 1/8
pada masing-masing hasil tersebut.  Kita menetapkan peluang 0 pada 16 hasil lainnya.  Dalam setiap kasus,
fungsi peluang itu disebut fungsi distribusi gabungan.
\end{example}

Kita merangkum gagasan-gagasan di atas dalam sebuah definisi.

\begin{definition}\label{def 4.4}
Misalkan $X_1, X_2, \ldots, X_n$ adalah peubah acak yang berkaitan dengan suatu percobaan.  Misalkan
ruang sampel (yaitu, himpunan hasil yang mungkin) bagi $X_i$ adalah himpunan $R_i$.  Maka
peubah acak gabungan\index{peubah acak!gabungan}\index{peubah acak gabungan} ${\bar
X} = (X_1, X_2,
\ldots, X_n)$ didefinisikan sebagai peubah acak yang hasilnya terdiri atas tupel-$n$ terurut,
dengan koordinat
ke-$i$ terletak dalam himpunan $R_i$.  Ruang sampel $\Omega$ bagi ${\bar X}$ adalah
hasil kali Kartesius dari himpunan-himpunan $R_i$:  
$$\Omega = R_1 \times R_2 \times \cdots \times R_n\ .$$
Fungsi distribusi gabungan\index{fungsi distribusi!gabungan}\index{fungsi distribusi
gabungan} bagi
${\bar X}$ adalah fungsi yang memberikan peluang bagi setiap hasil
${\bar X}$.
\end{definition}

\begin{example}(Lanjutan Contoh~\ref{exam 4.91}){\label{exam 4.93}}
Sekarang kita meninjau penetapan peluang dalam contoh di atas.  Untuk peubah acak
${\bar X}$, peluang sebarang hasil $(a_1, a_2, a_3)$ hanyalah hasil kali
peluang $P(X_i = a_i)$, untuk $i = 1, 2, 3$.  Namun, untuk
${\bar Y}$, peluang yang ditetapkan pada hasil $(1, 1, 0)$ bukanlah hasil kali
peluang $P(Y_1 = 1)$, $P(Y_2 = 1)$, dan $P(Y_3 = 0)$.  Perbedaan antara kedua
situasi ini adalah bahwa nilai $X_i$ tidak memengaruhi nilai $X_j$ jika $i \ne j$,
sedangkan nilai $Y_i$ dan $Y_j$ saling memengaruhi.  Sebagai contoh, jika $Y_1 = 1$, maka
$Y_2$ tidak mungkin sama dengan 0.  Hal ini mendorong definisi berikut.
\end{example}

\begin{definition}\label{def 4.4.5}
Peubah-peubah acak $X_1$,~$X_2$, \ldots,~$X_n$ disebut \emx {saling
bebas}\index{kebebasan peubah acak!saling}\index{peubah acak yang saling\\ 
bebas}\index{peubah acak!kebebasan saling} jika
\begin{eqnarray*}
&&P(X_1 = r_1, X_2 = r_2, \ldots, X_n = r_n)   \\
&& \;\;\;\;\;\;\;\;\;\;\;\;\;\;\; = P(X_1 = r_1) P(X_2 = r_2) \cdots  P(X_n = r_n)
\end{eqnarray*} 
untuk sebarang pilihan $r_1, r_2, \ldots, r_n$.  Jadi, jika $X_1,~X_2, \ldots,~X_n$ saling
bebas, fungsi distribusi gabungan peubah acak 
$${\bar X} = (X_1, X_2, \ldots, X_n)$$
hanyalah hasil kali fungsi-fungsi distribusi
masing-masing.  Jika dua peubah\linebreak[4] acak saling bebas, secara lebih
singkat kita katakan keduanya \emx {be-\linebreak[4] bas.}\index{kebebasan peubah\\
acak}\index{peubah acak!kebebasan}
\end{definition}

\begin{example}\label{exam 4.94} Dalam sekelompok 60 orang, banyaknya orang yang merokok atau tidak merokok
dan menderita atau tidak menderita kanker dilaporkan seperti dalam Tabel~\ref{table 4.1}.
\begin{table}
\centering
\begin{tabular}{l|cc|c}
             &\hspace{.15in}Tidak merokok\hspace{.15in}&\hspace{.15in}Merokok\hspace{.15in}&\hspace{.15in}Total\hspace{.15in}  \\ \hline
 Tanpa kanker  & 40             & 10             & 50 \\
 Kanker      & \hspace{.1in}7 &\hspace{.075in}3& 10 \\ \hline
 Total      & 47          & 13     & 60 \\
\end{tabular}
\caption{Merokok dan kanker.}
\label{table 4.1}
\end{table}
Misalkan $\Omega$ adalah ruang sampel yang terdiri atas 60 orang ini.  Seseorang dipilih secara acak
dari kelompok tersebut.  Misalkan $C(\omega) = 1$ jika orang ini menderita kanker dan~0 jika tidak, serta $S(\omega) =
1$ jika orang ini merokok dan~0 jika tidak. Maka distribusi gabungan $\{C,S\}$ diberikan dalam 
Tabel~\ref{table 4.2}.
\begin{table}
\centering
\begin{tabular}{lr|lcr}
  &    &            &  S      \\
  &    & \hspace{.15in} 0&  & 1\hspace{.15in}\\ \hline
  & 0  &  40/60     &  &10/60  \\ 
C &    &            &  &       \\
  & 1  & \hspace{.1in}7/60  &  &3/60   \\
\end{tabular}
\caption{Distribusi gabungan.}
\label{table 4.2}
\end{table}
\noindent Sebagai contoh, $P(C = 0, S = 0) = 40/60$, $P(C = 0, S = 1) = 10/60$, dan seterusnya.  Distribusi
peubah acak masing-masing disebut \emx {distribusi
marginal.}\index{fungsi distribusi marginal}\index{fungsi distribusi!marginal}  Distribusi marginal
$C$ dan $S$ adalah:
$$ p_C = \pmatrix{ 0     & 1     \cr 50/60 & 10/60 \cr},
$$

$$ p_S = \pmatrix{ 0     & 1     \cr 47/60 & 13/60 \cr}.
$$
Peubah acak $S$ dan $C$ tidak bebas karena
\begin{eqnarray*}
 P(C = 1,S = 1)   &=& \frac 3{60} = .05\ , \\
P(C = 1)P(S = 1) &=& \frac {10}{60} \cdot \frac
{13}{60} = .036\ .
\end{eqnarray*}
Perhatikan bahwa kita juga dapat melihatnya dari fakta bahwa
\begin{eqnarray*} P(C = 1|S = 1) &=& \frac 3{13} = .23\ , \\
P(C = 1) &=& \frac 16 = .167\ .
\end{eqnarray*}

\end{example}

\subsection*{Proses Percobaan Bebas}

Kajian peubah acak dilanjutkan dengan mempertimbangkan kelas-kelas khusus peubah acak. 
Salah satu kelas yang akan kita pelajari adalah kelas \emx {percobaan bebas.}

\begin{definition}\label{def 5.5} Suatu barisan peubah acak $X_1$,~$X_2$, \dots,~$X_n$
yang saling bebas dan memiliki distribusi yang sama disebut barisan
percobaan bebas atau \emx {proses percobaan bebas.}\index{proses percobaan bebas}

Proses percobaan bebas muncul secara alami dengan cara berikut.  Kita memiliki satu
percobaan dengan ruang sampel $R = \{r_1,r_2,\dots,r_s\}$ dan fungsi distribusi
$$m_X = \pmatrix{ r_1 & r_2 & \cdots & r_s \cr p_1 & p_2 & \cdots & p_s\cr}\ .
$$

Kita mengulangi percobaan ini sebanyak $n$ kali.  Untuk menggambarkan percobaan keseluruhan, sebagai ruang
sampel kita memilih ruang
$$
\Omega = R \times R \times\cdots\times R,
$$ yang terdiri atas semua barisan yang mungkin $\omega = (\omega_1,\omega_2,\dots,\omega_n)$ dengan
nilai setiap $\omega_j$ dipilih dari $R$.  Kita menetapkan fungsi distribusi berupa
\emx {distribusi hasil kali}
$$ m(\omega) = m(\omega_1)\cdot\ \ldots\ \cdot m(\omega_n)\ ,
$$ 
dengan $m(\omega_j) = p_k$ ketika $\omega_j = r_k$.  Kemudian kita menggunakan $X_j$ untuk menyatakan koordinat ke-$j$
dari hasil $(r_1, r_2, \ldots, r_n)$.  Peubah-peubah acak
$X_1$,~\dots,~$X_n$ membentuk suatu proses percobaan bebas.
\end{definition}

\begin{example}\label{exam 5.6.1} Suatu percobaan terdiri atas pelemparan sebuah dadu tiga kali.  Misalkan
$X_i$ menyatakan hasil lemparan ke-$i$, untuk $i = 1, 2, 3$.  Fungsi distribusi yang
sama untuk ketiganya adalah
$$ m_i = \pmatrix{ 1 & 2 & 3 & 4 & 5 & 6 \cr 1/6 & 1/6 & 1/6 & 1/6 & 1/6 & 1/6 \cr}.
$$

Ruang sampelnya adalah $R^3 = R \times R \times R$ dengan $R = \{1,2,3,4,5,6\}$.  Jika
$\omega = (1,3,6)$, maka $X_1(\omega) = 1$, $X_2(\omega) = 3$, dan
$X_3(\omega) = 6$, yang menunjukkan bahwa lemparan pertama menghasilkan 1, lemparan kedua 3, dan lemparan ketiga
6.  Peluang yang ditetapkan pada sebarang titik sampel adalah
$$ m(\omega) = \frac16 \cdot \frac16 \cdot \frac16 = \frac1{216}\ .
$$
\end{example}

\begin{example}\label{exam 5.7} Selanjutnya, perhatikan suatu proses percobaan Bernoulli dengan peluang $p$
untuk sukses pada setiap percobaan.  Misalkan $X_j(\omega) = 1$ jika hasil ke-$j$ adalah sukses dan
$X_j(\omega) = 0$ jika gagal.  Maka $X_1$,~$X_2$, \dots,~$X_n$ adalah suatu proses
percobaan bebas.  Setiap $X_j$ memiliki fungsi distribusi yang sama,
$$ m_j = \pmatrix{ 0 & 1 \cr q & p  \cr},
$$ dengan $q = 1 - p$.

Jika $S_n = X_1 + X_2 +\cdots + X_n$, maka
$$ P(S_n = j) = {n \choose j} p^{j} q^{n - j}\ ,  
$$ dan distribusi $S_n$ adalah distribusi binomial $b(n,p,j)$.
\end{example}

\subsection*{Rumus Bayes}     

Dalam contoh-contoh kita, kita telah membahas peluang bersyarat dengan bentuk
berikut: Dengan mengetahui hasil tahap kedua dari suatu percobaan dua tahap, tentukan
peluang suatu hasil pada tahap pertama.  Kita telah menyebutkan bahwa peluang-peluang
ini disebut \emx {peluang Bayes.}

Sekarang kita kembali ke perhitungan peluang Bayes yang lebih umum.  Misalkan
kita memiliki himpunan kejadian $H_1,$~$H_2$, \dots,~$H_m$ yang saling lepas berpasangan
dan sedemikian sehingga ruang sampel $\Omega$ memenuhi persamaan
$$
\Omega = H_1 \cup H_2 \cup\cdots\cup H_m\ .
$$
Kejadian-kejadian ini kita sebut \emx {hipotesis.}\index{hipotesis}  Kita juga memiliki kejadian $E$ yang memberi kita
sejumlah informasi tentang hipotesis mana yang benar.  Kejadian ini kita sebut \emx {bukti.}

Sebelum menerima bukti, kita memiliki sekumpulan \emx {peluang
prior}\index{peluang prior} $P(H_1)$,                \hfill\break
~$P(H_2)$, \dots,~$P(H_m)$ bagi hipotesis-hipotesis tersebut. 
Jika kita mengetahui hipotesis yang benar, kita mengetahui peluang bukti itu.  Dengan
kata lain, kita mengetahui $P(E|H_i)$ untuk semua $i$.  Kita ingin menentukan peluang
hipotesis dengan syarat bukti telah diperoleh.  Artinya, kita ingin menentukan peluang
bersyarat $P(H_i|E)$.  Peluang-peluang ini disebut \emx {peluang
posterior.}\index{peluang posterior}

Untuk menentukan peluang-peluang ini, kita menuliskannya dalam bentuk

\begin{equation}
P(H_i|E) = \frac{P(H_i \cap E)}{P(E)}\ . \label{eq 4.1}
\end{equation}
Kita dapat menghitung pembilang dari informasi yang diberikan melalui
\begin{equation}
P(H_i \cap E) = P(H_i)P(E|H_i)\ . \label{eq 4.2}  
\end{equation}
Karena tepat satu dari kejadian $H_1$,~$H_2$, \dots,~$H_m$ dapat terjadi, kita
dapat menuliskan peluang $E$ sebagai

$$
P(E) = P(H_1 \cap E)  + P(H_2 \cap E)  + \cdots + P(H_m \cap E)\ .
$$
Dengan menggunakan Persamaan~\ref{eq 4.2}, bentuk di atas dapat dilihat sama dengan

\begin{equation}
   P(H_1)P(E|H_1) + P(H_2)P(E|H_2) + \cdots + P(H_m)P(E|H_m) \ .
\label{eq 4.3}                                               
\end{equation}
Dengan menggunakan (\ref{eq 4.1}), (\ref{eq 4.2}), dan (\ref{eq 4.3}),
kita memperoleh \emx {rumus Bayes}:\index{rumus Bayes}
$$
P(H_i|E) = \frac{P(H_i)P(E|H_i)}{\sum_{k = 1}^m P(H_k)P(E|H_k)}\ .
$$

Meskipun rumus ini sangat terkenal, kita jarang akan menggunakannya.  Jika banyaknya
hipotesis sedikit, perhitungan ukuran pohon yang sederhana mudah
dilakukan, seperti dalam contoh-contoh kita.  Jika banyaknya hipotesis besar, kita sebaiknya
menggunakan komputer.
\par
Peluang Bayes sangat sesuai untuk diagnosis medis.  Seorang
dokter ingin mengetahui penyakit mana di antara beberapa kemungkinan yang mungkin diderita pasien.  Dokter tersebut
mengumpulkan bukti berupa hasil sejumlah tes.  Dari
kajian statistik, dokter dapat menentukan peluang prior berbagai
penyakit sebelum tes dilakukan, serta peluang hasil tes tertentu
dengan syarat pasien menderita penyakit tertentu.  Yang ingin diketahui dokter adalah peluang
posterior penyakit tersebut dengan syarat hasil-hasil tes telah diketahui.

\begin{example}\label{exam 4.10}
Seorang dokter sedang menentukan apakah seorang pasien menderita salah satu dari tiga penyakit,
$d_1$,~$d_2$, atau~$d_3$.  Dua tes akan dilakukan, yang masing-masing memberikan hasil
positif $(+)$ atau negatif $(-)$.  Ada empat pola tes yang mungkin,
$+{}+$,~$+{}-$, $-{}+$, dan~$-{}-$.  Catatan nasional menunjukkan
bahwa untuk 10{,}000 orang yang menderita salah satu dari ketiga penyakit ini, distribusi
penyakit dan hasil tes adalah seperti dalam Tabel~\ref{table 4.3}.

%\Mplace{\Table \paste{maintab}4 \endTable}
\begin{table}
\centering
\begin{tabular}{c|c|cccc} 
             & Jumlah penderita & \multicolumn{4}{c}{\underbar{Hasil tes}}  \\
  Penyakit    & penyakit ini\hspace{.25in} &\hspace{.2in}+\hspace{.1in}+\hspace{.2in}
&\hspace{.1in}+\hspace{.1in}--\hspace{.2in}&\hspace{.2in}--\hspace{.1in}+\hspace{.2in}
&\hspace{.1in}--\hspace{.1in}--\hspace{.2in}\\ \hline
  $ d_{1} $  &\hspace{.075in}3215&2110            &\hspace{.15in}301&\hspace{.1in}704  & \hspace{.1in}100\\
  $ d_{2} $  &\hspace{.075in}2125&\hspace{.1in}396&\hspace{.15in}132&1187  & \hspace{.1in}410\\
  $ d_{3} $  &\hspace{.075in}4660&\hspace{.1in}510&\hspace{.075in}3568&\hspace{.15in}73  & \hspace{.1in}509\\ \hline
Total        &10000  &&&&\\ \hline
\end{tabular}
\caption{Data penyakit.}
\label{table 4.3}
\end{table}

Dari data ini, kita dapat menaksir peluang prior masing-masing
penyakit dan, dengan syarat suatu penyakit tertentu, peluang suatu hasil tes
tertentu.  Sebagai contoh, peluang prior penyakit $d_1$ dapat ditaksir
sebesar $3215/10{,}000 = .3215$.  Peluang hasil tes $+{}-$ dengan syarat
penyakit $d_1$ dapat ditaksir sebesar $301/3215 = .094$.
\par
Sekarang kita dapat menggunakan rumus Bayes untuk menghitung berbagai peluang posterior.  Program
komputer {\bf Bayes}\index{Bayes (program)} menghitung peluang-peluang posterior ini.  Hasil
untuk contoh ini ditampilkan dalam Tabel~\ref{table 4.35}.


\begin{table}
\centering
\begin{tabular}{lccc}
                                    &  $d_1$ &  $d_2$ &  $d_3$  \\
+\hspace{.1in}+                     & .700   & .131   & .169    \\
+\hspace{.125in}--                    & .075   & .033   & .892    \\
\hspace{.025in}--\hspace{.125in}+     &
.358   & .604   & .038    \\
\hspace{.025in}--\hspace{.15in}--  &
.098   & .403   & .499    \\
\end{tabular}
\caption{Peluang posterior.}
\label{table 4.35}
\end{table}
Dari hasil tersebut kita mencatat bahwa ketika hasil tes adalah $++$, penyakit $d_1$
memiliki peluang yang jauh lebih tinggi daripada dua penyakit lainnya.  Ketika hasilnya $+-$,
hal ini berlaku untuk penyakit $d_3$.  Ketika hasilnya $-+$, hal ini berlaku untuk penyakit $d_2$.
Perhatikan bahwa pernyataan-pernyataan ini mungkin dapat diduga dengan melihat datanya. Jika hasilnya $--$, 
penyebab yang paling mungkin adalah $d_3$, tetapi peluang bahwa pasien menderita $d_2$ hanya sedikit lebih kecil.
Jika kita melihat data dalam kasus ini, kita dapat melihat bahwa mungkin sulit untuk menduga mana di antara kedua
penyakit $d_2$ dan $d_3$ yang lebih mungkin.
\end{example}
  
Contoh terakhir kita menunjukkan bahwa kita harus berhati-hati ketika peluang prior
kecil.

\begin{example}\label{exam 4.11}
Seorang dokter memberikan tes untuk suatu kanker tertentu\index{kanker} kepada seorang pasien.  Sebelum hasil
tes diperoleh, satu-satunya bukti yang dimiliki dokter adalah bahwa 1 dari~1000 perempuan menderita
kanker ini.  Pengalaman menunjukkan bahwa dalam 99 persen kasus ketika
kanker ada, hasil tes positif; dan dalam 95 persen kasus ketika
kanker tidak ada, hasil tes negatif.  Jika ternyata hasil tes positif,
peluang berapakah yang seharusnya dokter tetapkan pada kejadian bahwa kanker ada? 
Bentuk lain pertanyaan ini adalah menanyakan frekuensi relatif positif 
palsu dan kanker.

Kita diberi bahwa %${\rm prior(cancer)} = .001$ and
$\mbox{prior(kanker)} = .001$ dan
$\mbox{prior(tidak\ kanker)} = .999$.  Kita juga mengetahui bahwa $P(+| \mbox{kanker}) =
.99$, $P(-|\mbox{kanker}) = .01$, $P(+|\mbox{tidak\ kanker}) = .05$, dan
$P(-|\mbox{tidak\ kanker}) = .95$.  Penggunaan data ini memberikan hasil yang ditampilkan dalam
Gambar~\ref{fig 4.5}.

\putfig{5truein}{PSfig4-5}{Diagram pohon maju dan terbalik.}{fig 4.5}

Sekarang kita melihat bahwa peluang adanya kanker dengan syarat hasil tes positif hanya
meningkat dari .001 menjadi .019.  Meskipun peningkatan ini hampir dua puluh kali lipat,
peluang bahwa pasien menderita kanker masih kecil.  Dengan kata lain,
di antara hasil positif, 98.1 persen merupakan positif palsu dan 1.9 persen merupakan 
kanker.  Ketika pertanyaan ini diajukan kepada sekelompok mahasiswa kedokteran tahun kedua, lebih dari separuh
mahasiswa keliru menduga bahwa peluangnya lebih besar dari .5.  
\end{example}

\subsection*{Catatan Sejarah}

Peluang bersyarat telah digunakan jauh sebelum didefinisikan secara formal.  Pascal dan Fermat
membahas \emx {masalah pembagian taruhan}:\index{masalah pembagian taruhan} dengan syarat tim A telah memenangi $m$
permainan dan tim B telah memenangi $n$ permainan, berapakah peluang bahwa A akan memenangi rangkaian tersebut?  (Lihat
Latihan~\ref{exer 5.1.11}--\ref{exer 5.1.13}.)  Ini jelas merupakan masalah peluang
bersyarat.  
\par
Dalam bukunya, Huygens\index{HUYGENS, C.} memberikan masalah berikut. Sebuah guci memuat
12 bola, terdiri atas 4 bola putih dan 8 bola hitam. Tiga pemain, A, B, dan C, mengambil bola
secara bergiliran dalam urutan tersebut tanpa melihat; pemain pertama yang memperoleh bola putih
menang. Tentukan rasio peluang menang A, B, dan C.\footnote{Ringkasan berdasarkan F.~N. David,
\emx {Games, Gods and Gambling} (London: Griffin, 1962), p.~119; ungkapan terjemahan modern
tidak diterjemahkan atau direproduksi.}
\par
Dari jawabannya, jelas bahwa Huygens bermaksud setiap bola dikembalikan setelah
diambil.  Namun, John Hudde\index{HUDDE, J.}, wali kota Amsterdam, mengasumsikan bahwa Huygens bermaksud melakukan
pengambilan sampel tanpa pengembalian dan berkorespondensi dengan Huygens tentang perbedaan
jawaban mereka.  Menurut ringkasan Hacking, masing-masing tidak berhasil memahami model yang
digunakan pihak lain.\index{HACKING, I.}\footnote{I.~Hacking, \emx {The Emergence of Probability}
(Cambridge: Cambridge University Press, 1975), p.~99. Ringkasan editorial; ungkapan sumber
tidak diterjemahkan atau direproduksi.}
\par
Pada masa terbitnya buku de~Moivre\index{de MOIVRE, A.}, \emx {The Doctrine of Chances,} perbedaan-perbedaan ini
telah dipahami dengan baik.  De~Moivre mendefinisikan kebebasan dan
kebergantungan sebagai berikut:
\begin{quote}
\indent Dua Kejadian bebas jika tidak memiliki hubungan satu dengan yang
lain, dan terjadinya salah satu tidak mempercepat ataupun menghambat
terjadinya yang lain.
\par
Dua Kejadian bergantung jika keduanya terhubung sedemikian rupa sehingga
Peluang terjadinya salah satu berubah akibat terjadinya kejadian yang
lain.\footnote{A.~de~Moivre, \emx {The Doctrine of Chances,} 3rd~ed. (New York:
Chelsea, 1967), p.~6.}
\end{quote}

De~Moivre menggunakan pengambilan sampel dengan dan tanpa pengembalian untuk menunjukkan bahwa
peluang terjadinya kedua kejadian bebas adalah hasil kali peluang
keduanya, dan bahwa untuk kejadian bergantung:
\begin{quote}
Peluang terjadinya dua Kejadian yang bergantung adalah hasil kali
Peluang terjadinya salah satu kejadian dengan Peluang terjadinya kejadian yang lain
ketika kejadian pertama dianggap telah terjadi; dan
Aturan yang sama berlaku untuk terjadinya sebanyak apa pun Kejadian yang dapat
ditentukan.\footnote{ibid, p.~7.}
\end{quote}
\par
Rumus yang kita sebut rumus Bayes, serta gagasan menghitung
peluang suatu hipotesis dengan syarat bukti tertentu, berasal dari sebuah esai terkenal karya
Thomas Bayes\index{BAYES, T.}.  Bayes adalah pendeta tertahbis di Tunbridge Wells dekat London. 
Minat matematikanya membuatnya terpilih sebagai anggota Royal Society pada~1742,
tetapi tidak satu pun hasilnya diterbitkan semasa hidupnya.  Karya yang
melandasi ketenarannya, ``An Essay Toward Solving a Problem in the Doctrine of
Chances," diterbitkan pada~1763, tiga tahun setelah
kematiannya.\footnote{T.~Bayes, ``An Essay Toward Solving a Problem in the Doctrine
of Chances," \emx {Phil.\ Trans.\ Royal Soc.\ London,} vol.~53 (1763),
pp.~370--418.}  Bayes meninjau beberapa konsep dasar peluang, lalu
membahas jenis baru masalah peluang balikan yang memerlukan penggunaan
peluang bersyarat.
\par
Dalam kajiannya tentang proses yang kini kita sebut percobaan Bernoulli, Bernoulli\index{BERNOULLI, J.} telah
membuktikan hukum bilangan besar yang terkenal, yang akan kita pelajari dalam Bab~\ref{chp
8}.  Teorema ini menjamin bahwa jika pelaku percobaan mengetahui peluang $p$
untuk sukses, ia dapat meramalkan bahwa proporsi sukses akan mendekati
nilai tersebut ketika jumlah percobaan diperbesar.  Bernoulli sendiri
menyadari bahwa dalam kebanyakan kasus yang menarik, nilai $p$ tidak diketahui, dan
memandang teoremanya sebagai langkah penting untuk menunjukkan bahwa $p$ dapat ditentukan melalui
percobaan.
\par
Untuk mengkaji masalah ini lebih lanjut, Bayes mulai dengan mengasumsikan bahwa peluang
$p$ untuk sukses itu sendiri ditentukan oleh suatu percobaan acak.  Ia bahkan
mengasumsikan bahwa percobaan tersebut membuat nilai $p$ memiliki peluang yang sama untuk
menjadi sebarang nilai di antara 0 dan 1.\choice{\footnote{Bayes menggunakan peluang
kontinu seperti yang dibahas dalam Bab 2 buku lengkap Grinstead-Snell.}}{}  Tanpa mengetahui nilai ini, kita
melakukan
$n$ percobaan dan mengamati $m$ sukses.  Bayes mengajukan masalah untuk menentukan
peluang bersyarat bahwa peluang $p$ yang tidak diketahui terletak di antara $a$
dan $b$.  Ia memperoleh jawaban:
$$
P(a \leq p < b | m {\mbox{\,\,sukses\,\, dalam}}\,\,n \,\,{\mbox{percobaan}}) 
= \frac {\int_a^b x^m(1 - x)^{n - m}\,dx}{\int_0^1 x^m(1 - x)^{n - m}\,dx}\ .
$$
\par
\choice{}{Dalam bagian berikutnya kita akan melihat bagaimana hasil ini diperoleh.}  Bayes jelas
ingin menunjukkan bahwa fungsi distribusi bersyarat, dengan syarat hasil dari
semakin banyak percobaan, menjadi terkonsentrasi di sekitar nilai $p$ yang sebenarnya. 
Jadi, Bayes berusaha menyelesaikan suatu \emx {masalah balikan.}  Perhitungan
integral-integral tersebut terlalu sulit untuk diselesaikan secara eksak kecuali untuk nilai
$j$ dan~$n$ yang kecil, sehingga Bayes mencoba metode pendekatan.  Metodenya tidak
begitu memuaskan, dan ada dugaan bahwa hal ini membuatnya enggan
menerbitkan hasil-hasilnya.
\par
Namun, makalahnya merupakan yang pertama dalam serangkaian kajian penting yang dilakukan
oleh Laplace, Gauss, dan matematikawan besar lainnya untuk menyelesaikan masalah balikan. 
Mereka mengkaji masalah ini dalam kaitannya dengan galat pengukuran dalam astronomi.  Jika
seorang astronom mengetahui nilai sebenarnya suatu jarak dan sifat
galat acak yang disebabkan oleh alat ukurnya, ia dapat meramalkan sifat probabilistik
pengukurannya.  Namun, pada kenyataannya ia menghadapi masalah balikan:
ia mengetahui sifat galat acak dan nilai-nilai pengukuran, lalu ingin
membuat inferensi tentang nilai sebenarnya yang tidak diketahui.
\par
Seperti dinyatakan Maistrov, rumus yang kita sebut rumus Bayes tidak
muncul dalam esai Bayes.  Laplace memberinya nama tersebut ketika ia mengkaji masalah-masalah
balikan ini.\index{MAISTROV, L.}\footnote{L.~E. Maistrov, \emx {Probability Theory: A Historical
Sketch,} trans.~and ed.~Samual Kotz (New York: Academic Press, 1974), p.~100.}
Perhitungan peluang balikan bersifat mendasar bagi statistika dan telah
melahirkan cabang penting statistika yang disebut analisis Bayesian, sehingga
esai singkat Bayes memberinya ketenaran abadi.

\exercises
\begin{LJSItem}

\i\label{exer 4.1.1} Anggaplah $E$ dan $F$ adalah dua kejadian dengan peluang
positif.  Tunjukkan bahwa jika $P(E|F) = P(E)$, maka $P(F|E) = P(F)$.

\i\label{exer 4.1.2} Sebuah koin dilempar tiga kali.  Berapakah peluang tepat dua
kepala muncul, dengan syarat bahwa
\begin{enumerate}
\item hasil pertama adalah kepala?

\item hasil pertama adalah ekor?

\item dua hasil pertama adalah kepala?

\item dua hasil pertama adalah ekor?

\item hasil pertama adalah kepala dan hasil ketiga adalah kepala?
\end{enumerate}

\i\label{exer 4.1.3} Sebuah dadu dilempar dua kali.  Berapakah peluang bahwa jumlah mata dadu
lebih besar dari 7, dengan syarat bahwa
\begin{enumerate}
\item hasil pertama adalah 4?

\item hasil pertama lebih besar dari 3?

\item hasil pertama adalah 1?

\item hasil pertama kurang dari 5?
\end{enumerate}

\i\label{exer 4.1.4} Sebuah kartu diambil secara acak dari setumpuk kartu.  Berapakah peluang
bahwa
\begin{enumerate}
\item kartu itu berhati, dengan syarat kartu itu merah?

\item nilainya lebih tinggi dari 10, dengan syarat kartu itu berhati? (Tafsirkan J,~Q, K,~A
sebagai 11,~12, 13,~14.)

\item kartu itu jack, dengan syarat kartu itu merah?
\end{enumerate} 

\i\label{exer 4.1.5} Sebuah koin dilempar tiga kali.  Tinjaulah kejadian-kejadian berikut.\newline
$A$: Kepala pada lemparan pertama.\newline
$B$: Ekor pada lemparan kedua.\newline
$C$: Kepala pada lemparan ketiga.\newline
$D$: Ketiga hasil sama (HHH atau TTT).\newline
$E$: Tepat satu kepala muncul.
\begin{enumerate}
\item Manakah di antara pasangan kejadian berikut yang saling bebas?\newline
(1) $A$, $B$\newline
(2) $A$, $D$\newline
(3) $A$, $E$\newline
(4) $D$, $E$

\item Manakah di antara rangkap tiga kejadian berikut yang saling
bebas?\newline
(1) $A$, $B$, $C$\newline
(2) $A$, $B$, $D$\newline
(3) $C$, $D$, $E$
\end{enumerate}

\i\label{exer 4.1.6} Dari setumpuk lima kartu yang masing-masing bernomor 2,~4, 6, 8, dan~10,
sebuah kartu diambil secara acak lalu dikembalikan.  Proses ini dilakukan tiga kali.  Berapakah
peluang bahwa kartu bernomor 2 terambil tepat dua kali, dengan syarat
jumlah nomor pada ketiga pengambilan adalah~12?

\i\label{exer 4.1.7} Sebuah koin dilempar dua kali.  Tinjaulah
kejadian-kejadian berikut.\newline
$A$: Kepala pada lemparan pertama.\newline
$B$: Kepala pada lemparan kedua.\newline
$C$: Kedua lemparan memberikan hasil yang sama.
\begin{enumerate}
\item Tunjukkan bahwa $A$,~$B$,~$C$ bebas berpasangan, tetapi tidak
saling bebas.

\item Tunjukkan bahwa $C$ bebas terhadap $A$ dan $B$, tetapi tidak terhadap $A \cap B$.
\end{enumerate}

\i\label{exer 4.1.8} Misalkan $\Omega = \{a,b,c,d,e,f\}$.  Anggaplah $m(a) = m(b) = 1/8$ dan
$m(c) = m(d) = m(e) = m(f) = 3/16$.  Misalkan $A$,~$B$, dan~$C$ adalah kejadian $A = \{d,e,a\}$, 
$B = \{c,e,a\}$, $C = \{c,d,a\}$.  Tunjukkan bahwa $P(A \cap B \cap C) = P(A)P(B)P(C)$, tetapi tidak ada dua dari
kejadian-kejadian ini yang saling bebas.
 
\i\label{exer 4.1.9} Berapakah peluang bahwa suatu keluarga dengan dua anak mempunyai
\begin{enumerate}
\item dua anak laki-laki, dengan syarat sekurang-kurangnya satu anaknya laki-laki?

\item dua anak laki-laki, dengan syarat anak pertamanya laki-laki?
\end{enumerate}

\i\label{exer 4.1.9.5}  Dalam Contoh~\ref{exam 4.1.5}, kita menggunakan Tabel Hayat (lihat Lampiran~C)
untuk menghitung suatu peluang bersyarat.  Bilangan 93{,}753 dalam tabel itu, yang bersesuaian dengan
laki-laki berusia 40 tahun, berarti bahwa dari seluruh laki-laki yang lahir di Amerika Serikat pada 1950, sebanyak 93.753\%
masih hidup pada 1990.  Apakah wajar menggunakan nilai ini sebagai taksiran peluang bahwa seorang laki-laki yang lahir
tahun ini akan bertahan hidup sampai usia 40 tahun? 

\i\label{exer 4.1.9.6} Simulasikan masalah Monty Hall.  Nyatakan dengan cermat setiap asumsi
yang Anda buat ketika menulis program.  Menurut Anda, versi masalah manakah yang
Anda simulasikan?

\i\label{exer 4.1.10} Dalam Contoh~\ref{exam 4.11}, seberapa besarkah peluang prior
kanker agar diperoleh peluang posterior .5 untuk kanker dengan syarat hasil tes
positif?

\i\label{exer 4.1.11} Dua kartu diambil dari setumpuk kartu bridge.  Berapakah peluang bahwa
kartu kedua yang diambil berwarna merah?

\i\label{exer 4.1.12} Jika $P(\tilde B) = 1/4$ dan $P(A|B) = 1/2$, berapakah $P(A \cap B)$?

\i\label{exer 4.1.13}
\begin{enumerate} 
\item Berapakah peluang bahwa pasangan bridge Anda memiliki tepat dua kartu as,
dengan syarat ia memiliki sekurang-kurangnya satu kartu as?

\item Berapakah peluang bahwa pasangan bridge Anda memiliki tepat dua kartu as,
dengan syarat ia memiliki kartu as sekop?
\end{enumerate}

\i\label{exer 4.1.14} Buktikan bahwa untuk sebarang tiga kejadian $A$,~$B$,~$C$ yang masing-masing memiliki
peluang positif dan memenuhi $P(A \cap B) > 0$,
$$
P(A \cap B \cap C) = P(A)P(B|A)P(C|A \cap B)\ .
$$

\i\label{exer 4.1.15} Buktikan bahwa jika $A$ dan $B$ saling bebas, maka pasangan berikut juga saling bebas:
\begin{enumerate}
\item $A$ dan $\tilde B$.

\item $\tilde A$ dan $\tilde B$.
\end{enumerate}

\i\label{exer 4.1.16} Seorang dokter menganggap bahwa seorang pasien menderita salah satu dari tiga penyakit $d_1$,~$d_2$,
atau~$d_3$.  Sebelum melakukan tes apa pun, ia menganggap setiap penyakit memiliki peluang yang sama. 
Ia melakukan suatu tes yang akan memberikan hasil positif dengan peluang .8 jika pasien
menderita $d_1$, .6 jika pasien menderita penyakit $d_2$, dan .4 jika pasien menderita penyakit $d_3$.  Dengan syarat
hasil tes positif, peluang berapakah yang sekarang seharusnya ditetapkan dokter
bagi ketiga penyakit yang mungkin itu?

\i\label{exer 4.1.17} Dalam suatu permainan poker, John memiliki kombinasi kartu yang sangat kuat dan bertaruh 5 dolar.  Peluang
bahwa Mary memiliki kombinasi yang lebih baik adalah .04.  Jika Mary memiliki kombinasi yang lebih baik, ia
akan menaikkan taruhan dengan peluang .9, tetapi dengan kombinasi yang lebih buruk ia hanya akan menaikkan taruhan
dengan peluang .1.  Jika Mary menaikkan taruhan, berapakah peluang bahwa ia memiliki
kombinasi yang lebih baik daripada John?

\i\label{exer 4.1.18} Model guci Polya\index{model guci Polya} untuk penularan adalah sebagai berikut: Kita
mulai dengan sebuah guci yang berisi satu bola putih dan satu bola hitam.  Pada setiap
detik kita memilih sebuah bola secara acak dari guci, mengembalikan bola itu, lalu menambahkan satu bola lagi
dengan warna yang terpilih.  Tulislah program untuk menyimulasikan model ini, dan lihat apakah
Anda dapat membuat prediksi tentang proporsi bola putih dalam guci
setelah sejumlah besar pengambilan.  Apakah dalam jangka panjang terdapat kecenderungan bahwa sebagian besar
bola memiliki warna yang sama?

\i\label{exer 4.1.19} Kita ingin mencari peluang bahwa dalam suatu pembagian kartu bridge setiap pemain
menerima satu kartu as.  Seorang mahasiswa berargumen sebagai berikut.  Tidak penting kepada siapa
kartu as pertama dibagikan.  Kartu as kedua harus dibagikan kepada salah satu dari tiga pemain lainnya dan
hal ini terjadi dengan peluang 3/4.  Kemudian kartu as berikutnya harus dibagikan kepada salah satu dari dua pemain, suatu
kejadian dengan peluang 1/2, dan akhirnya kartu as terakhir harus dibagikan kepada pemain yang
belum memiliki kartu as.  Hal ini terjadi dengan peluang 1/4.  Peluang bahwa
semua kejadian ini terjadi adalah hasil kali $(3/4)(1/2)(1/4) = 3/32$.  Apakah
argumen ini benar?

\i\label{exer 4.1.20} Satu koin dalam kumpulan 65 koin mempunyai dua sisi kepala.  Koin-koin lainnya seimbang.  Jika sebuah
koin yang dipilih secara acak dari kumpulan itu lalu dilempar menghasilkan kepala 6 kali
berturut-turut, berapakah peluang bahwa koin tersebut adalah koin berkepala dua?

\i\label{exer 4.1.21} Anda diberi dua guci dan lima puluh bola.  Separuh bola berwarna putih dan
separuhnya hitam.  Anda diminta membagikan bola-bola tersebut ke dalam kedua guci tanpa
pembatasan jumlah bola dari masing-masing warna dalam suatu guci.  Bagaimanakah Anda harus
membagikan bola-bola itu agar peluang memperoleh bola
putih menjadi maksimum jika sebuah guci dipilih secara acak dan sebuah bola diambil secara acak? 
Berikan alasan bagi jawaban Anda.

\i\label{exer 4.1.22} Sebuah koin seimbang dilempar $n$ kali.  Tunjukkan bahwa peluang bersyarat
munculnya kepala pada suatu lemparan tertentu, dengan syarat terdapat total $k$ kepala dalam $n$
lemparan, adalah $k/n$ $(k > 0)$.

\i\label{exer 4.1.23} (Johnsonbough\index{JOHNSONBOUGH, R.}\footnote{R. Johnsonbough, ``Problem
\#103," \emx {Two Year College Math Journal,} vol.~8 (1977), p.~292.})  Sebuah koin dengan peluang
$p$ untuk kepala dilempar $n$ kali.  Misalkan $E$ adalah kejadian ``kepala diperoleh pada
lemparan pertama" dan $F_k$ adalah kejadian ``tepat $k$ kepala diperoleh."  Untuk pasangan
$(n,k)$ manakah $E$ dan $F_k$ saling bebas?

\i\label{exer 4.1.24} Misalkan $A$ dan $B$ adalah kejadian sedemikian sehingga $P(A|B) = P(B|A)$, $P(A
\cup B) = 1$, dan $P(A \cap B) > 0$.  Buktikan bahwa $P(A) > 1/2$.

\i\label{exer 4.1.25} (Chung\index{CHUNG, K. L.}\footnote{K. L. Chung, \emx {Elementary
Probability Theory With Stochastic Processes, 3rd ed.} (New York:  Springer-Verlag, 1979), p.~152.})
Di London, hujan turun pada separuh hari.  Peramal cuaca benar sebanyak 2/3 dari seluruh kesempatan;
artinya, peluang hujan turun dengan syarat ia meramalkan hujan dan peluang hujan
tidak turun dengan syarat ia meramalkan tidak akan hujan, keduanya sama dengan 2/3.  Ketika hujan
diramalkan, Tn.\ Pickwick\index{Pickwick, Tn.} membawa payungnya.  Ketika hujan tidak diramalkan, ia
membawanya dengan peluang 1/3.  Carilah
\begin{enumerate}
\item peluang bahwa Pickwick tidak membawa payung, dengan syarat hujan turun.

\item peluang bahwa ia membawa payung, dengan syarat hujan tidak turun.
\end{enumerate}

\i\label{exer 4.1.26} Teori peluang digunakan dalam suatu perkara pengadilan
terkenal: \emx {People v. Collins.}\index{Collins, People v.}\index{People v. Collins}\footnote{M.~W.
Gray, ``Statistics and the Law," \emx {Mathematics Magazine,} vol.~56 (1983), pp.~67--81.}  Dalam
perkara ini, sebuah dompet dirampas dari seorang lansia di pinggiran Los Angeles.  Sepasang orang yang terlihat
berlari dari tempat kejadian digambarkan sebagai seorang pria kulit hitam berjanggut\index{janggut} dan
berkumis\index{kumis}, serta seorang perempuan berambut pirang yang dikuncir\index{kuncir}.  Para saksi mengatakan
mereka melarikan diri dengan mobil yang sebagian berwarna kuning.  Malcolm dan Janet Collins ditangkap.  Malcolm berkulit hitam dan,
meskipun bercukur bersih ketika ditangkap, menunjukkan bukti bahwa belum lama sebelumnya ia berjanggut
dan berkumis.  Janet berambut pirang dan biasanya menguncir rambutnya.  Mereka
mengendarai Lincoln yang sebagian berwarna kuning.  Jaksa menghadirkan seorang profesor matematika
sebagai saksi; ia mengusulkan bahwa sekumpulan peluang konservatif untuk
ciri-ciri yang dicatat para saksi adalah seperti yang ditampilkan dalam Tabel~\ref{table 4.4}.
\begin{table}[h]
\centering
\begin{tabular}{lc}
{\rm pria berkumis}    & 1/4          \cr
{\rm perempuan berambut pirang} & 1/3          \cr
{\rm perempuan berambut dikuncir}   & 1/10         \cr
{\rm pria kulit hitam berjanggut} & 1/10          \cr
{\rm pasangan antar-ras dalam mobil} & 1/1000 \cr
{\rm mobil yang sebagian berwarna kuning}    & 1/10    \cr
\end{tabular}
\caption{Peluang-peluang dalam perkara Collins.}
\label{table 4.4}
\end{table}
\par
Jaksa kemudian berargumen bahwa peluang semua
ciri ini dipenuhi oleh pasangan yang dipilih secara acak adalah hasil kali
peluang-peluangnya, yaitu 1/12{,}000{,}000, yang sangat kecil.  Ia menyatakan bahwa hal ini merupakan bukti
tanpa keraguan yang beralasan bahwa para terdakwa bersalah.  Juri menyetujuinya
dan menjatuhkan putusan bersalah atas perampokan tingkat dua.

Jika Anda menjadi pengacara pasangan Collins, bagaimanakah Anda akan membantah
argumen di atas?  (Banding dalam perkara ini dibahas dalam 
Latihan~\ref{sec 5.1}.\ref{exer 9.2.23}.)

\i\label{exer 4.1.27} Seorang mahasiswa mendaftar ke Harvard dan Dartmouth.  Ia menaksir bahwa ia
memiliki peluang .5 untuk diterima di Dartmouth dan .3 untuk diterima di
Harvard.  Ia juga menaksir bahwa peluang ia diterima di keduanya
adalah .2.  Berapakah peluang bahwa ia diterima di Dartmouth jika ia
diterima di Harvard?  Apakah kejadian ``diterima di Harvard" bebas terhadap
kejadian ``diterima di Dartmouth"?

\i\label{exer 4.1.28} Luxco, produsen grosir bola lampu, memiliki dua pabrik.  Pabrik A
menjual bola lampu dalam lot yang masing-masing terdiri atas 1000 bola lampu biasa dan 2000 bola lampu \emx {softglow}.
Pengambilan sampel acak menunjukkan bahwa secara rata-rata cenderung terdapat
sekitar 2 bola lampu biasa yang rusak dan 11 bola lampu softglow yang rusak per lot.  Di pabrik B,
komposisi lot dibalik---terdapat 2000 bola lampu biasa dan 1000 bola lampu softglow per lot---dan
cenderung terdapat 5 bola lampu biasa yang rusak dan 6 bola lampu softglow yang rusak per lot.

Manajer pabrik A menyatakan, ``Jelas kami produsen yang lebih baik; tingkat
bola lampu rusak kami adalah .2 persen dan .55 persen, dibandingkan dengan .25 persen dan
.6 persen milik B.  Untuk bola lampu biasa maupun softglow, kami lebih baik sebesar setengah dari sepersepuluh
persen."

``Sebaliknya," bantah manajer B, ``setiap lot kami yang berisi 3000 bola lampu
hanya mengandung 11 bola lampu rusak, sedangkan lot A yang berisi 3000 bola lampu mengandung 13.  Jadi, tingkat
bola lampu rusak kami sebesar .37 persen mengalahkan tingkat mereka sebesar .43 persen."
\par
Siapakah yang benar?

\i\label{exer 4.1.29} Dengan menggunakan Tabel Hayat tahun 1981 yang diberikan dalam Lampiran~C, carilah
peluang bahwa seorang laki-laki yang berusia~60 tahun pada 1981 hidup sampai usia~80 tahun.  Carilah
peluang yang sama untuk seorang perempuan.

\i\label{exer 4.1.30}
\begin{enumerate}
\item Telah terjadi badai salju dan Helen berusaha berkendara dari Woodstock\index{Woodstock}
ke Tunbridge\index{Tunbridge}, yang terhubung seperti pada graf teratas dalam Gambar~\ref{fig
4.51}.  Di sini $p$ dan $q$ adalah peluang bahwa kedua jalan dapat dilalui.  Berapakah
peluang bahwa Helen dapat pergi dari Woodstock ke Tunbridge?

\item Sekarang misalkan Woodstock dan Tunbridge terhubung seperti pada graf tengah 
dalam Gambar~\ref{fig 4.51}.
Berapakah sekarang peluang bahwa Helen dapat pergi dari $W$ ke $T$?  Perhatikan bahwa jika kita
memandang jalan-jalan itu sebagai komponen suatu sistem, maka dalam
(a) dan (b) kita telah menghitung
\emx {keandalan}\index{keandalan suatu sistem} sistem yang komponennya tersusun (a)~\emx {secara
seri} dan (b)~\emx {secara paralel.}

\item Sekarang misalkan $W$ dan $T$ terhubung seperti pada graf terbawah dalam Gambar~\ref{fig 4.51}.
Carilah peluang bahwa Helen dapat pergi dari $W$ ke $T$.  \emx {Petunjuk}: Jika
jalan dari $C$ ke $D$ tidak dapat dilalui, jalan itu dapat dianggap tidak ada; jika jalan itu
dapat dilalui, carilah cara menggunakan bagian (b) dua kali.
\end{enumerate}

\putfig{3.5truein}{PSfig4}{Dari Woodstock ke Tunbridge.}{fig 4.51} %%3.5truein


\i\label{exer 4.1.31} Misalkan $A_1$,~$A_2$, dan $A_3$ adalah kejadian, dan misalkan $B_i$ menyatakan
$A_i$ atau komplemennya $\tilde A_i$.  Maka terdapat delapan pilihan yang mungkin untuk
rangkap tiga $(B_1, B_2, B_3)$.  Buktikan bahwa kejadian $A_1$,~$A_2$,~$A_3$ saling bebas jika dan
hanya jika $$P(B_1 \cap B_2 \cap B_3) = P(B_1)P(B_2)P(B_3)\ ,$$
untuk kedelapan pilihan yang mungkin bagi rangkap tiga $(B_1, B_2, B_3)$.

\i\label{exer 5.1.4} Empat perempuan, A,~B, C, dan~D, menitipkan topi mereka, lalu topi-topi tersebut dikembalikan
secara acak.  Misalkan $\Omega$ adalah himpunan semua permutasi yang mungkin dari A,~B, C,~D.  Misalkan
$X_j = 1$ jika perempuan ke-$j$ menerima kembali topinya sendiri dan 0 jika tidak.  Apakah
distribusi $X_j$?  Apakah semua $X_i$ saling bebas?

\i\label{exer 5.1.5} Sebuah kotak memuat bilangan 1 hingga 10.  Sebuah bilangan diambil secara acak.  Misalkan
$X_1$ adalah bilangan yang diambil.  Bilangan ini dikembalikan, lalu kesepuluh bilangan dicampur.  Bilangan kedua,
$X_2$, diambil.  Carilah distribusi $X_1$ dan $X_2$.  Apakah
$X_1$ dan $X_2$ saling bebas?  Jawablah pertanyaan yang sama jika bilangan pertama tidak dikembalikan
sebelum bilangan kedua diambil.

\i\label{exer 5.1.6} Sebuah dadu dilempar dua kali.  Misalkan $X_1$ dan $X_2$ menyatakan hasilnya. 
Definisikan $X = \min(X_1, X_2)$.  Carilah distribusi $X$.

\istar\label{exer 5.1.7} \begin{sloppypar} Dengan diketahui bahwa $P(X = a) = r$, $P(\max(X,Y) = a) = s$,
dan  
$P(\min(X,Y) = a) = t$, tunjukkan bahwa Anda dapat menentukan $u = P(Y = a)$ dalam bentuk $r$,~$s$,
dan~$t$.\end{sloppypar}

\i\label{exer 5.1.8} Sebuah koin seimbang dilempar tiga kali.  Misalkan $X$ adalah banyaknya kepala yang
muncul pada dua lemparan pertama dan $Y$ adalah banyaknya kepala yang muncul pada lemparan ketiga. 
Berikan distribusi
\begin{enumerate}
\item peubah acak $X$ dan $Y$.

\item peubah acak $Z = X + Y$.

\item peubah acak $W = X - Y$.
\end{enumerate}

\i\label{exer 5.1.10} Anggaplah peubah acak $X$ dan $Y$ memiliki distribusi gabungan
yang diberikan dalam Tabel~\ref{table 4.5}.
\begin{table}
\centering
\begin{tabular}{lr|rrrrr} 
     &     &  $Y$  &       &        &        \\
     &     &   -1  &  0    &  1     &  2     \\ \hline
$X$  & -1  &   0   & 1/36  &  1/6   & 1/12   \\
     &  0  & 1/18  &  0    &  1/18  &  0     \\
     &  1  &   0   & 1/36  &  1/6   & 1/12   \\
     &  2  & 1/12  &  0    &  1/12  & 1/6    \\
\end{tabular}
\caption{Distribusi gabungan.}
\label{table 4.5}
\end{table}

\begin{enumerate}
\item Berapakah $P(X \geq 1\ \mbox {dan\ } Y \leq 0)$?

\item Berapakah peluang bersyarat bahwa $Y \leq 0$ dengan syarat $X = 2$?

\item Apakah $X$ dan $Y$ saling bebas?

\item Bagaimanakah distribusi $Z = XY$?
\end{enumerate}

\i\label{exer 5.1.11} Dalam \emx {masalah pembagian taruhan}\index{masalah pembagian taruhan}, yang dibahas dalam 
catatan sejarah pada Bagian~\ref{sec 3.2}, dua pemain, A dan B, memainkan serangkaian poin dalam suatu
permainan; pemain A memenangi setiap poin dengan peluang $p$, sedangkan pemain B memenangi setiap poin dengan
peluang
$q = 1 - p$.  Pemain pertama yang memenangi $N$ poin memenangi permainan.  Anggaplah $N = 3$.  Misalkan $X$
adalah peubah acak yang bernilai~1 jika pemain A memenangi rangkaian dan 0 jika tidak.  Misalkan
$Y$ adalah peubah acak yang nilainya sama dengan banyaknya poin yang dimainkan dalam suatu permainan.  Carilah
distribusi $X$ dan $Y$ ketika $p = 1/2$.  Apakah $X$ dan $Y$ saling bebas dalam kasus ini?  Jawablah
pertanyaan yang sama untuk kasus $p = 2/3$.

\i\label{exer 5.1.12} Surat-menyurat antara Pascal\index{PASCAL, B.} dan Fermat\index{FERMAT, P.}, yang
sering dianggap telah memulai teori peluang, sebagian besar membahas \emx {masalah pembagian
taruhan} yang diuraikan dalam Latihan~\ref{exer 5.1.11}.  Pascal dan Fermat meninjau masalah
menentukan pembagian taruhan yang adil jika permainan harus dihentikan ketika pemain pertama telah memenangi
$r$ permainan dan pemain kedua telah memenangi $s$ permainan, dengan $r < N$ dan $s < N$.  Misalkan $P(r,s)$ adalah
peluang bahwa pemain A memenangi permainan jika ia telah memenangi $r$ poin dan pemain B telah memenangi
$s$ poin.  Maka
\begin{enumerate}
\item $P(r,N) = 0$ jika $r < N$,
\item $P(N,s) = 1$ jika $s < N$,
\item $P(r,s) = pP(r + 1,s) + qP(r,s + 1)$ jika $r < N$ dan $s < N$;
\end{enumerate} dan (1),~(2), serta~(3) menentukan $P(r,s)$ untuk $r \leq N$ dan $s \leq N$.  Pascal
menggunakan fakta-fakta ini untuk mencari $P(r,s)$ dengan bekerja mundur: Mula-mula ia memperoleh $P(N - 1,j)$ untuk $j =
N - 1$,~$N - 2$, \dots,~0; kemudian, dari nilai-nilai ini, ia memperoleh
$P(N - 2,j)$ untuk $j = N - 1$,~$N - 2$, \dots,~0 dan, dengan terus bekerja mundur, memperoleh semua
nilai $P(r,s)$.  Tulislah program untuk menghitung $P(r,s)$ bagi $N$,~$a$, $b$, dan~$p$ yang diberikan.  \emx {
Peringatan}: Ikutilah Pascal dan Anda akan dapat menjalankan $N = 100$; gunakan rekursi dan Anda
tidak akan dapat menjalankan $N = 20$.

\i\label{exer 5.1.13} Fermat menyelesaikan \emx {masalah pembagian taruhan} (lihat Latihan~\ref{exer
5.1.11}) sebagai berikut: Ia menyadari bahwa masalah tersebut sulit karena berbagai kemungkinan jalannya
permainan tidak memiliki peluang yang sama.  Sebagai contoh, ketika pemain pertama memerlukan dua kemenangan lagi
dan pemain kedua memerlukan tiga kemenangan lagi, dua kemungkinan jalannya rangkaian bagi pemain pertama
adalah WLW dan LWLW.  Kedua barisan ini tidak memiliki peluang yang sama.  Untuk menghindari kesulitan ini, Fermat
memperpanjang permainan dengan menambahkan permainan fiktif sehingga rangkaian berlangsung selama jumlah maksimum permainan yang
diperlukan (empat dalam kasus ini).  Ia memperoleh hasil-hasil dengan peluang yang sama dan, pada hakikatnya, menggunakan
segitiga Pascal untuk menghitung $P(r,s)$.  Tunjukkan bahwa cara ini menghasilkan suatu \emx {rumus} untuk $P(r,s)$
bahkan dalam kasus $p \ne 1/2$.

\i\label{exer 5.1.14} Yankees bertanding melawan Dodgers dalam suatu seri dunia.  Yankees memenangi
setiap pertandingan dengan peluang .6.  Berapakah peluang bahwa Yankees memenangi seri tersebut?  (Seri
dimenangkan oleh tim pertama yang memenangi empat pertandingan.)

\i\label{exer 5.1.15} C.~L. Anderson\index{ANDERSON, C. L.}\footnote{C.~L. Anderson, ``Note on the
Advantage of First Serve," \emx {Journal of Combinatorial Theory,} Series~A, vol.~23 (1977),
p.~363.} menggunakan argumen Fermat untuk \emx {masalah pembagian taruhan} guna membuktikan hasil berikut
yang berasal dari J.~G. Kingston\index{KINGSTON, J. G.}.  Anda memainkan \emx {permainan perolehan poin} (lihat
Latihan~\ref{exer 5.1.11}), tetapi pada setiap poin, ketika Anda melakukan servis, Anda menang dengan peluang $p$, dan
ketika lawan melakukan servis, Anda menang dengan peluang $\bar{p}$.  Anda melakukan servis terlebih dahulu, tetapi Anda
dapat memilih salah satu dari dua aturan servis berikut: pada aturan pertama, Anda
bergantian melakukan servis (tenis)\index{tenis}, sedangkan pada aturan kedua, orang yang melakukan servis terus melakukannya
hingga kehilangan satu poin, lalu pemain lain melakukan servis (raketbol)\index{raketbol}.  Pemain
pertama yang memenangi $N$ poin memenangi permainan.  Masalahnya adalah menunjukkan bahwa peluang
memenangi permainan sama pada kedua aturan.
\begin{enumerate}

\item Tunjukkan bahwa menurut aturan mana pun, Anda akan melakukan servis paling banyak untuk $N$ poin dan lawan Anda
paling banyak untuk $N - 1$ poin.

\item Perpanjang jumlah poin menjadi $2N - 1$ sehingga Anda melakukan servis untuk $N$ poin dan lawan Anda
melakukan servis untuk $N - 1$ poin.  Sebagai contoh, Anda melakukan servis untuk setiap poin tambahan yang diperlukan agar jumlah servis Anda menjadi $N$, lalu
lawan Anda melakukan servis untuk setiap poin tambahan yang diperlukan agar jumlah servisnya menjadi $N - 1$. 
Pemenangnya sekarang adalah orang yang memenangi poin terbanyak dalam permainan yang diperpanjang.  Tunjukkan bahwa
memainkan poin-poin tambahan ini tidak mengubah pemenang.

\item Tunjukkan bahwa (a) dan (b) membuktikan bahwa peluang Anda memenangi permainan sama menurut
kedua aturan.
\end{enumerate}

\i\label{exer 5.1.15.5}
Dalam soal sebelumnya, anggaplah $p = 1 - \bar{p}$.
\begin{enumerate}
\item  Tunjukkan bahwa menurut aturan servis mana pun, pemain pertama akan lebih sering menang daripada pemain kedua
jika dan hanya jika $p > .5$.
\item
Dalam bola voli,\index{bola voli} sebuah tim hanya dapat memenangi poin ketika sedang melakukan servis.  Dengan demikian, setiap
``rally" berakhir dengan pemberian poin kepada tim yang melakukan servis atau dengan perpindahan servis
ke tim lain.  Tim pertama yang memenangi $N$ poin memenangi permainan.  (Di sini kita mengabaikan 
syarat tambahan bahwa tim pemenang harus unggul sekurang-kurangnya dua poin pada akhir
permainan.)  Anggaplah setiap tim memiliki peluang yang sama untuk memenangi rally ketika melakukan servis, yakni 
$p = 1 - \bar{p}$.  Tunjukkan bahwa dalam kasus ini, tim yang melakukan servis pertama akan menang lebih dari separuh
kesempatan, selama $p > 0$.  (Jika $p = 0$, permainan tidak pernah berakhir.)  \emx {Petunjuk}: Definisikan $p'$ sebagai
peluang bahwa suatu tim memenangi poin berikutnya, dengan syarat tim itu sedang melakukan servis.  Jika kita tulis $q = 1 -
p$, maka dapat ditunjukkan bahwa 
$$p' = \frac p{1-q^2}\ .$$
Jika sekarang permainan ini ditinjau dengan cara yang sedikit berbeda, dapat dilihat bahwa aturan servis kedua
dalam soal sebelumnya dapat digunakan dengan mengganti $p$ oleh $p'$.
\end{enumerate} 

\i\label{exer 5.1.19} Suatu kombinasi poker terdiri atas 5 kartu yang dibagikan dari setumpuk 52 kartu.  Misalkan $X$
dan $Y$ berturut-turut adalah banyaknya kartu as dan raja dalam kombinasi poker.  Carilah distribusi
gabungan $X$ dan $Y$.

\i\label{exer 5.1.24.5} Misalkan $X_1$ dan $X_2$ adalah peubah acak yang saling bebas dan misalkan $Y_1 =
\phi_1(X_1)$ dan $Y_2 = \phi_2(X_2)$.
\begin{enumerate}

\item Tunjukkan bahwa
$$ P(Y_1 = r, Y_2 = s) = \sum_{\phi_1(a) = r \atop \phi_2(b) = s} P(X_1 = a, X_2 = b)\ .
$$

\item Dengan menggunakan (a), tunjukkan bahwa $P(Y_1 = r, Y_2 = s) = P(Y_1 = r)P(Y_2 = s)$ sehingga $Y_1$ dan
$Y_2$ saling bebas.
\end{enumerate}

\i\label{exer 4.1.32} Misalkan $\Omega$ adalah ruang sampel suatu percobaan.  Misalkan $E$ adalah
kejadian dengan $P(E) > 0$ dan definisikan $m_E(\omega)$ dengan
$m_E(\omega) = m(\omega|E)$.  Buktikan bahwa $m_E(\omega)$ adalah fungsi distribusi pada
$E$, yakni bahwa $m_E(\omega) \geq 0$ dan $\sum_{\omega\in\Omega}
m_E(\omega) = 1$.  Fungsi $m_E$ disebut \emx {distribusi bersyarat dengan syarat $E$.}

\i\label{exer 4.1.33} Anda diberi dua guci yang masing-masing berisi dua koin bias.  Koin-koin dalam
guci~I menghasilkan kepala dengan peluang~$p_1$, sedangkan koin-koin dalam guci~II menghasilkan
kepala dengan peluang $p_2 \ne p_1$.  Anda diberi pilihan (a)~memilih
sebuah guci secara acak dan melempar kedua koin di dalam guci itu atau (b)~memilih satu koin
dari setiap guci dan melempar kedua koin tersebut.  Anda memenangi hadiah jika kedua koin menghasilkan
kepala.  Tunjukkan bahwa pilihan (a) lebih menguntungkan.

\i\label{exer 4.1.34} Buktikan bahwa jika $A_1$,~$A_2$, \dots,~$A_n$ adalah kejadian-kejadian saling bebas
yang didefinisikan pada ruang sampel $\Omega$ dan jika $0 < P(A_j) < 1$ untuk semua $j$, maka
$\Omega$ harus memiliki sekurang-kurangnya $2^n$ titik.

\i\label{exer 4.1.35} Buktikan bahwa jika
$$
P(A|C) \geq P(B|C) \mbox{\,\,dan\,\,} P(A|\tilde C) \geq P(B|\tilde C)\ ,
$$
maka $P(A) \geq P(B)$.

\i\label{exer 4.1.36} Sebuah koin berada dalam salah satu dari $n$ kotak.  Peluang bahwa koin itu berada dalam kotak ke-$i$
adalah $p_i$.  Jika Anda mencari dalam kotak ke-$i$ dan koin itu ada di sana, Anda menemukannya dengan
peluang $a_i$.  Tunjukkan bahwa peluang~$p$ bahwa koin berada dalam kotak ke-$j$,
dengan syarat Anda telah mencari dalam kotak ke-$i$ dan tidak menemukannya, adalah
$$
p =  \left \{ \matrix{
                 p_j/(1-a_ip_i),&\,\,\, \mbox{jika} \,\,\, j \ne i,\cr
                 (1 - a_i)p_i/(1 - a_ip_i),&\,\,\,\mbox{jika} \,\, j = i.\cr}\right. 
$$

\i\label{exer 4.1.37} George Wolford\index{WOLFORD, G.} mengusulkan variasi berikut dari
masalah Linda (lihat Latihan~\ref{sec 1.2}.\ref{exer 1.2.24}).  
Registrar membawa kartu registrasi John dan Mary, lalu menjatuhkannya ke dalam
genangan air.  Ketika memungutnya, ia tidak dapat membaca nama-nama pada kartu, tetapi pada kartu pertama
yang dipungutnya ia dapat mengenali Mathematics 23 dan Government 35, sedangkan pada kartu kedua
ia hanya dapat mengenali Mathematics 23.  Ia bertanya apakah Anda dapat membantunya
menentukan kartu mana yang milik Mary.  Anda tahu bahwa Mary menyukai ilmu pemerintahan, tetapi
tidak menyukai matematika.  Anda tidak mengetahui apa pun tentang John dan menganggapnya sebagai
mahasiswa Dartmouth pada umumnya.  Berdasarkan hal ini, Anda menaksir:
\vskip .1in
$$\begin{array}{ll}
P(\mbox {Mary\ mengambil\ Government\ 35})  &= .5\ , \\
P(\mbox {Mary\ mengambil\ Mathematics\ 23}) &= .1\ , \\
P(\mbox {John\ mengambil\ Government\ 35})  &= .3\ , \\
P(\mbox {John\ mengambil\ Mathematics\ 23}) &= .2\ .
\end{array}
$$
\vskip .1in
Anggaplah pilihan mata kuliah mereka merupakan kejadian-kejadian saling bebas.  Tunjukkan bahwa
kartu yang menampilkan Mathematics 23 dan Government 35 lebih mungkin merupakan milik Mary
daripada milik John.  Kekeliruan konjungsi yang disebutkan dalam masalah Linda adalah
menganggap bahwa kejadian ``Mary mengambil Mathematics 23 dan Government 35"
lebih mungkin daripada kejadian ``Mary mengambil Mathematics 23."  Mengapa kita tidak
melakukan kekeliruan tersebut di sini?

\i\label{exer 4.1.38} (Diusulkan oleh Eisenberg dan Ghosh\index{EISENBERG,
B.}\index{GHOSH, B. K.}\footnote{B. Eisenberg and B.~K. Ghosh, ``Independent Events in a Discrete
Uniform Probability Space," \emx {The American Statistician,} vol.~41, no.~1 (1987),
pp.~52--56.})  Setumpuk kartu permainan dapat dinyatakan sebagai hasil kali Cartesius
$$
\mbox{Tumpukan} =  \mbox{Kembang} \times \mbox{Nilai}\ ,
$$
dengan $\mbox{Kembang} = \{\clubsuit,\diamondsuit,\heartsuit,\spadesuit\}$ dan %
$\mbox{Nilai} = \{2,3,\dots,10,{\mbox J},{\mbox Q},{\mbox K},{\mbox A}\}$.  Hal ini
hanya berarti bahwa setiap kartu dapat dipandang sebagai pasangan terurut seperti
$(\diamondsuit,2)$.  Yang dimaksud dengan \emx {kejadian kembang}\index{kejadian kembang} adalah sebarang kejadian $A$ di dalam
Tumpukan yang hanya dinyatakan menurut kembang.  Sebagai contoh, jika $A$ adalah ``kembang
berwarna merah," maka
$$
A = \{\diamondsuit,\heartsuit\} \times \mbox{Nilai}\ ,
$$
sehingga $A$ terdiri atas semua kartu berbentuk $(\diamondsuit,r)$ atau
$(\heartsuit,r)$, dengan $r$ suatu nilai kartu sembarang.  Demikian pula, \emx {kejadian nilai kartu}\index{kejadian nilai kartu} adalah sebarang
kejadian yang hanya dinyatakan menurut nilai kartu.
\begin{enumerate}
\item Tunjukkan bahwa jika $A$ adalah sebarang kejadian kembang dan $B$ adalah sebarang kejadian nilai kartu, maka $A$
dan $B$ \emx {saling bebas.}  (Kita dapat menyatakannya secara ringkas dengan mengatakan bahwa
kembang dan nilai kartu saling bebas.)

\item Buanglah kartu as sekop.  Tunjukkan bahwa sekarang tidak ada kejadian kembang nontrivial (yakni,
bukan himpunan kosong maupun seluruh ruang) $A$ yang bebas terhadap sebarang
kejadian nilai kartu nontrivial $B$.  \emx {Petunjuk}: Di sini kebebasan tereduksi menjadi
$$
c/51 = (a/51) \cdot (b/51)\ ,
$$
dengan $a$,~$b$,~$c$ berturut-turut adalah ukuran $A$,~$B$, dan~$A \cap B$.  Akibatnya,
51 harus membagi $ab$, sehingga 3 harus membagi salah satu dari $a$ dan $b$,
sedangkan 17 membagi yang lain.  Akan tetapi, kemungkinan ukuran kejadian kembang dan nilai kartu tidak memungkinkan
hal ini.

\item Tunjukkan bahwa tumpukan kartu pada (b) bagaimanapun memang memiliki
pasangan kejadian nontrivial yang saling bebas, $A$,~$B$.  \emx {Petunjuk}: Carilah 2 kejadian
$A$ dan $B$ yang masing-masing berukuran 3 dan 17 serta beririsan di satu titik.

\item Tambahkan satu kartu joker ke setumpuk kartu lengkap.  Tunjukkan bahwa sekarang tidak ada pasangan $A$,~$B$
berupa kejadian nontrivial yang saling bebas.  \emx {Petunjuk}: Lihat petunjuk pada
(b); 53 adalah bilangan prima.
\end{enumerate}
\par
\medskip
\noindent
Soal-soal berikut diusulkan oleh Stanley Gudder\index{GUDDER, S.} dalam artikelnya
``Do Good Hands Attract?"\footnote{S. Gudder, ``Do Good Hands Attract?" \emx {
Mathematics Magazine,} vol.~54, no.~1 (1981), pp.~13--16.}  Ia mengatakan bahwa
kejadian $A$ \emx {menarik}\index{kejadian!daya tarik} kejadian $B$ jika $P(B|A) > P(B)$ dan \emx {
menolak}\index{kejadian!daya tolak}
$B$ jika $P(B|A) < P(B)$.

\i\label{exer 4.1.39} Misalkan $R_i$ adalah kejadian bahwa pemain ke-$i$ dalam permainan poker memiliki
royal flush.  Tunjukkan bahwa satu royal flush (A,K,Q,J,10 dari satu kembang) menarik royal flush lainnya, 
yakni $P(R_2|R_1) > P(R_2)$.  Tunjukkan bahwa royal flush menolak full house.

\i\label{exer 4.1.40} Buktikan bahwa $A$ menarik $B$ jika dan hanya jika $B$ menarik $A$.  Oleh karena itu, kita
dapat mengatakan bahwa $A$ dan $B$ \emx {saling menarik} jika $A$ menarik $B$.

\i\label{exer 4.1.41} Buktikan bahwa $A$ tidak menarik maupun menolak $B$ jika dan hanya jika $A$ dan
$B$ saling bebas.

\i\label{exer 4.1.42} Buktikan bahwa $A$ dan $B$ saling menarik jika dan hanya jika $P(B|A) >
P(B|\tilde A)$.

\i\label{exer 4.1.43} Buktikan bahwa jika $A$ menarik $B$, maka $A$ menolak $\tilde B$.

\i\label{exer 4.1.44} Buktikan bahwa jika $A$ menarik $B$ maupun $C$, dan $A$ menolak $B \cap C$,
maka $A$ menarik $B \cup C$.  Adakah contoh ketika $A$ menarik
$B$ maupun $C$, tetapi menolak $B \cup C$?

\i\label{exer 4.1.45} Buktikan bahwa jika $B_1$,~$B_2$, \dots,~$B_n$ saling lepas dan
secara bersama-sama mencakup seluruh ruang, dan jika $A$ menarik suatu $B_i$, maka $A$ harus menolak
suatu $B_j$.

\i\label{exer 4.1.46} 
\begin{enumerate}
\item
Misalkan Anda mencari di dalam meja tulis Anda sebuah surat dari beberapa waktu lalu.  Meja tulis itu memiliki delapan
laci, dan Anda menaksir bahwa peluang surat berada dalam setiap laci tertentu adalah 10\% (jadi terdapat
peluang 20\% bahwa surat sama sekali tidak berada dalam meja).  Sekarang misalkan Anda mulai mencari 
secara sistematis di meja itu, satu laci demi satu.  Selain itu, misalkan Anda belum menemukan
surat dalam $i$ laci pertama, dengan $0 \le i \le 7$.  Misalkan $p_i$ menyatakan peluang bahwa
surat akan ditemukan dalam laci berikutnya, dan misalkan $q_i$ menyatakan peluang bahwa surat 
akan ditemukan dalam suatu laci setelahnya (baik $p_i$ maupun $q_i$ merupakan peluang bersyarat karena
keduanya didasarkan pada asumsi bahwa surat tidak berada dalam $i$ laci pertama).  Tunjukkan bahwa
nilai-nilai $p_i$ meningkat dan nilai-nilai $q_i$ menurun.  (Soal ini berasal dari Falk dkk.\footnote{R.\
Falk, A.\ Lipson, and C.\ Konold, ``The ups and downs of the hope function in a fruitless search,"
in \emx{Subjective Probability,} G.\ Wright and P.\ Ayton, (eds.) (Chichester: Wiley, 1994), pgs.
353-377.})\index{FALK, R.}\index{LIPSON, A.}\index{KONOLD, C.}
\item
Data berikut dimuat dalam sebuah artikel di Wall Street Journal.\footnote{C. Crossen, ``Fright by the
numbers:  Alarming disease data are frequently flawed," \emx{Wall Street Journal,} 11 April 1996, p.
B1.}\index{Wall Street Journal}\index{CROSSEN, C.} Pada usia 20, 30, 40, 50, dan 60 tahun, peluang seorang 
perempuan di A.S.\ terkena kanker dalam sepuluh tahun berikutnya berturut-turut adalah 0.5\%, 1.2\%, 3.2\%, 6.4\%, dan 10.8\%.
Pada rangkaian usia yang sama, peluang seorang perempuan di A.S.\ pada akhirnya terkena kanker
berturut-turut adalah 39.6\%, 39.5\%, 39.1\%, 37.5\%, dan 34.2\%.  Menurut Anda, apakah soal pada bagian (a) memberikan
penjelasan bagi data ini? 
\end{enumerate}

\i\label{exer 4.1.47} 
Berikut ini dua variasi masalah Monty Hall yang dibahas oleh Granberg.\footnote{D.
Granberg, ``To switch or not to switch," in
\emx{The power of logical thinking,} M. vos~Savant, (New York: St.\ Martin's
1996).}\index{GRANBERG, D.}\index{masalah Monty Hall}
\begin{enumerate}
\item
Misalkan semuanya sama, kecuali bahwa Monty lupa mencari tahu terlebih dahulu pintu mana yang menyembunyikan mobil.
Dengan semangat ``pertunjukan harus terus berjalan," ia menebak pintu mana di antara kedua pintu yang akan dibuka dan
beruntung membuka pintu yang menyembunyikan seekor kambing.  Sekarang, haruskah peserta berpindah pilihan?
\item
Anda telah lama mengamati pertunjukan tersebut dan mendapati bahwa mobil diletakkan di balik pintu A dalam 45\% kesempatan,
di balik pintu B dalam 40\% kesempatan, dan di balik pintu C dalam 15\% kesempatan. Anggaplah segala hal lain
tentang pertunjukan itu tetap sama.  Sekali lagi Anda memilih pintu A. Monty membuka pintu yang menyembunyikan kambing dan
menawarkan kesempatan kepada Anda untuk berpindah pilihan. Haruskah Anda berpindah?  Misalkan Anda sebelumnya sudah mengetahui bahwa Monty akan memberi Anda
kesempatan untuk berpindah.  Apakah seharusnya sejak awal Anda memilih pintu A?
\end{enumerate}

\end{LJSItem}

%At one point, the following \choice command starting giving errors,
%so I commented it out (I also commented out the ending }).
%\choice{}{\section{Continuous Conditional Probability}\label{sec 4.2}
\section{Peluang Bersyarat Kontinu}\label{sec 4.2}

Dalam situasi ketika ruang sampel bersifat kontinu, kita akan mengikuti prosedur yang sama seperti pada
bagian sebelumnya.  Jadi, sebagai contoh, jika $X$ adalah peubah acak kontinu dengan fungsi densitas
$f(x)$, dan jika $E$ adalah kejadian dengan peluang positif, kita mendefinisikan fungsi
densitas bersyarat\index{densitas bersyarat}\index{fungsi densitas!bersyarat} dengan rumus
$$ f(x|E) = \left \{ \matrix{
 f(x)/P(E), &  \mbox{jika} \,\,x \in E,       \cr
         0, & \mbox{jika}\,\,x \not \in E.  \cr}\right.        
$$
Maka, untuk sebarang kejadian $F$, kita mempunyai
$$ 
P(F|E) = \int_F f(x|E)\,dx\ .
$$
Ekspresi $P(F|E)$ disebut peluang bersyarat $F$ dengan syarat $E$.  Seperti pada
bagian sebelumnya, kita dapat dengan mudah memperoleh ekspresi alternatif bagi peluang ini:
$$
P(F|E) = \int_F f(x|E)\,dx = \int_{E\cap F} \frac {f(x)}{P(E)}\,dx = \frac {P(E\cap F)}{P(E)}\ .
$$
\par
Kita dapat memandang fungsi densitas bersyarat sebagai fungsi yang bernilai 0 kecuali pada $E$, dan
dinormalisasi agar integralnya di atas $E$ sama dengan 1.  Perhatikan bahwa jika densitas semula adalah
densitas seragam yang bersesuaian dengan suatu percobaan ketika semua kejadian berukuran sama
\emx {berpeluang sama,} maka hal yang sama berlaku bagi densitas
bersyarat.

\begin{example}\label{exam 4.12}
Dalam percobaan pemutar\index{pemutar} (bdk.\ Contoh~\ref{exam 2.1.1}),
misalkan kita mengetahui bahwa pemutar berhenti dengan ujung penunjuk pada separuh bagian atas
lingkaran, $0 \leq x \leq 1/2$.  Berapakah peluang bahwa $1/6 \leq x
\leq 1/3$?

Di sini $E = [0,1/2]$, $F = [1/6,1/3]$, dan $F \cap E = F$.  Oleh karena itu,
\begin{eqnarray*}
P(F|E) &=& \frac {P(F \cap E)}{P(E)} \\
       &=& \frac {1/6}{1/2} \\
       &=& \frac 13\ ,
\end{eqnarray*}
yang masuk akal karena ukuran $F$ adalah 1/3 ukuran $E$.  Fungsi densitas
bersyarat di sini diberikan oleh

$$
f(x|E) =  \left \{ \matrix{
               2, & \mbox{jika}\,\,\, 0   \leq x < 1/2, \cr
               0, & \mbox{jika}\,\,\, 1/2 \leq x < 1.\cr}\right.
$$
Jadi, fungsi densitas bersyarat hanya tidak nol pada $[0,1/2]$, dan bersifat
seragam di sana.
\end{example}

\begin{example}\label{exam 4.13}
Dalam permainan anak panah\index{anak panah} (bdk.\ Contoh~\ref{exam 2.2.2}),
misalkan kita mengetahui bahwa anak panah jatuh pada separuh bagian atas sasaran.  Berapakah
peluang bahwa jaraknya dari pusat kurang dari 1/2?

Di sini $E = \{\,(x,y) : y \geq 0\,\}$ dan $F = \{\,(x,y) : x^2 + y^2 <
(1/2)^2\,\}$.  Oleh karena itu,
\begin{eqnarray*}
P(F|E) & = & \frac {P(F \cap E)}{P(E)} = \frac {(1/\pi)[(1/2)(\pi/4)]}
{(1/\pi)(\pi/2)} \\
       & = & 1/4\ .
\end{eqnarray*}
Di sini sekali lagi, ukuran $F \cap E$ adalah 1/4 ukuran $E$.  Fungsi densitas
bersyaratnya adalah
$$
f((x,y)|E) =  \left \{ \matrix{
        f(x,y)/P(E) = 2/\pi, &\mbox{jika}\,\,\,(x,y) \in E, \cr
        0,                   &\mbox{jika}\,\,\,(x,y) \not \in E.\cr}\right.
$$
\end{example}

\begin{example}\label{exam 4.14}
Kita kembali ke densitas eksponensial\index{densitas eksponensial}\index{fungsi
densitas!eksponensial} (bdk.\ Contoh~\ref{exam 2.2.7.5}).  Misalkan kita mengamati sebongkah
plutonium-239.  Percobaan kita terdiri atas menunggu suatu emisi, lalu menyalakan pencatat waktu, dan
mencatat lamanya waktu $X$ yang berlalu sampai emisi berikutnya.   Pengalaman menunjukkan bahwa
$X$ memiliki densitas eksponensial dengan suatu parameter $\lambda$, yang bergantung pada ukuran
bongkahan.  Misalkan ketika kita melakukan percobaan ini, kita melihat bahwa pencatat waktu menunjukkan $r$ detik
dan masih berjalan.  Berapakah peluang bahwa tidak terjadi emisi selama tambahan $s$ detik?

Misalkan $G(t)$ adalah peluang bahwa partikel berikutnya dipancarkan setelah waktu $t$.  Maka
\begin{eqnarray*}
G(t) & = & \int_t^\infty \lambda e^{-\lambda x}\,dx \\
     & = & \left.-e^{-\lambda x}\right|_t^\infty = e^{-\lambda t}\ .
\end{eqnarray*}

Misalkan $E$ adalah kejadian ``partikel berikutnya dipancarkan setelah waktu $r$" dan $F$ adalah kejadian ``partikel
berikutnya dipancarkan setelah waktu $r + s$."  Maka
\begin{eqnarray*}
P(F|E) & = & \frac {P(F \cap E)}{P(E)} \\
       & = & \frac {G(r + s)}{G(r)} \\
       & = & \frac {e^{-\lambda(r + s)}}{e^{-\lambda r}} \\
       & = & e^{-\lambda s}\ .
\end{eqnarray*}


Hal ini memberi tahu kita fakta yang cukup mengejutkan bahwa peluang kita harus
menunggu tambahan $s$ detik untuk suatu emisi, dengan syarat belum terjadi emisi selama $r$ detik,
\emx {tidak bergantung} pada waktu $r$.  Sifat ini (disebut \emx {
tanpa ingatan}\index{sifat tanpa ingatan}) diperkenalkan dalam Contoh~\ref{exam 2.2.7.5}. 
Ketika mencoba memodelkan berbagai fenomena, sifat ini membantu kita memutuskan apakah
densitas eksponensial sesuai.  
\par
Fakta bahwa densitas eksponensial bersifat tanpa ingatan berarti bahwa wajar untuk menganggap
bahwa jika seseorang menjumpai sebongkah isotop radioaktif pada suatu waktu acak, maka lama waktu
sampai emisi berikutnya memiliki densitas eksponensial dengan parameter yang sama seperti waktu
antaremisi.
Sebuah contoh terkenal, yang disebut ``paradoks bus,"\index{paradoks bus}, mengganti
emisi dengan bus.  Paradoks yang tampak ini muncul dari dua fakta berikut:  1) Jika Anda mengetahui bahwa
secara rata-rata bus datang setiap 30 menit, maka jika Anda tiba di halte bus pada waktu acak,
secara rata-rata Anda seharusnya hanya perlu menunggu 15 menit untuk sebuah bus, dan 2) Karena waktu
kedatangan bus dimodelkan dengan densitas eksponensial, kapan pun Anda tiba,
secara rata-rata Anda harus menunggu 30 menit untuk sebuah bus.
\par
Pembaca sekarang dapat melihat bahwa dalam Latihan~\ref{sec 2.2}.\ref{exer 2.2.8.5}, 
\ref{sec 2.2}.\ref{exer 2.2.8.6}, dan \ref{sec 2.2}.\ref{exer
2.2.11}, kita meminta simulasi peluang bersyarat berdasarkan berbagai asumsi
tentang distribusi waktu antarkedatangan.  Jika seseorang membuat asumsi yang wajar tentang
distribusi ini, seperti asumsi dalam Latihan~\ref{sec 2.2}.\ref{exer 2.2.8.6}, maka waktu 
tunggu rata-rata lebih mendekati separuh waktu antarkedatangan rata-rata.
\end{example}

\subsection*{Kejadian Bebas}

Jika $E$ dan $F$ adalah dua kejadian dengan peluang positif dalam ruang sampel kontinu, maka seperti pada
ruang sampel diskret, kita mendefinisikan $E$ dan $F$ sebagai \emx {
bebas}\index{kejadian!bebas}\index{kebebasan kejadian} jika 
$P(E|F) = P(E)$ dan $P(F|E) = P(F)$.  Seperti sebelumnya, masing-masing persamaan di atas mengakibatkan
persamaan yang lain, sehingga untuk melihat apakah dua kejadian saling bebas, hanya satu dari persamaan ini yang perlu
diperiksa.  Selain itu, jika $E$ dan $F$ saling bebas, maka $P(E \cap F) = P(E)P(F)$.

\begin{example}(Lanjutan Contoh~\ref{exam 4.12}){\label{exam 4.15}}
Dalam permainan anak panah\index{anak panah} (lihat Contoh~\ref{exam 4.12}), misalkan $E$ adalah kejadian bahwa
anak panah jatuh pada separuh bagian \emx {atas} sasaran ($y \geq 0$) dan $F$ adalah
kejadian bahwa anak panah jatuh pada separuh bagian \emx {kanan} sasaran ($x \geq
0$).  Maka $P(E \cap F)$ adalah peluang bahwa anak panah berada dalam kuadran pertama
sasaran, dan
\begin{eqnarray*}
P(E \cap F) & = & \frac 1\pi \int_{E \cap F} 1\,dxdy \\
            & = & \mbox{Luas}\,(E\cap F)         \\
            & = & \mbox{Luas}\,(E)\,\mbox{Luas}\,(F)  \\
            & = & \left(\frac 1\pi \int_E 1\,dxdy\right) \left(\frac 1\pi \int_F
1\,dxdy\right) \\
            & = & P(E)P(F)
\end{eqnarray*}
sehingga $E$ dan $F$ saling bebas.  Hal yang membuat perhitungan ini berhasil adalah bahwa kejadian
$E$ dan $F$ dinyatakan dengan membatasi koordinat yang berbeda.  Gagasan ini dibuat 
lebih tepat\linebreak[4] di bawah.
\end{example}

\subsection*{Fungsi Densitas dan Distribusi Kumulatif Gabungan}

Dengan cara yang serupa dengan peubah acak diskret, kita dapat mendefinisikan fungsi densitas
gabungan dan fungsi distribusi kumulatif untuk peubah acak kontinu multidimensi.  

\begin{definition}\label{def 4.5}
Misalkan $X_1,~X_2, \ldots,~X_n$ adalah peubah acak kontinu yang berkaitan dengan suatu percobaan, dan
misalkan ${\bar X} = (X_1,~X_2, \ldots,~X_n)$.  Maka fungsi distribusi kumulatif
gabungan\index{fungsi distribusi kumulatif!gabungan}\index{fungsi distribusi kumulatif 
\\ gabungan} dari
${\bar X}$ didefinisikan oleh
$$F(x_1, x_2, \ldots, x_n) = P(X_1 \le x_1, X_2 \le x_2, \ldots, X_n \le x_n)\ .$$
Fungsi densitas gabungan\index{fungsi densitas gabungan}\index{fungsi densitas!gabungan} dari
${\bar X}$ memenuhi persamaan berikut:
$$F(x_1, x_2, \ldots, x_n) = \int_{-\infty}^{x_1} \int_{-\infty}^{x_2} \cdots
\int_{-\infty}^{x_n} f(t_1, t_2, \ldots t_n)\,dt_ndt_{n-1}\ldots dt_1\ .$$
\end{definition}

Dengan notasi di atas, dapat ditunjukkan secara langsung bahwa
\begin{equation}
f(x_1, x_2, \ldots, x_n) = {{\partial^n F(x_1, x_2, \ldots, x_n)}\over
{\partial x_1 \partial x_2 \cdots \partial x_n}}\ .\label{eq 4.4}
\end{equation}

\subsection*{Peubah Acak Bebas}

Seperti pada peubah acak diskret, kita dapat mendefinisikan kebebasan bersama untuk
peubah acak kontinu.

\begin{definition}\label{def 4.6} Misalkan $X_1$,~$X_2$, \ldots,~$X_n$ adalah peubah acak
kontinu dengan fungsi distribusi kumulatif $F_1(x),~F_2(x), \ldots,~F_n(x)$.  Maka
peubah-peubah acak ini \emx {saling bebas}\index{kebebasan peubah acak!bersama} jika
$$F(x_1, x_2, \ldots, x_n) = F_1(x_1)F_2(x_2) \cdots  F_n(x_n)$$   
untuk sebarang pilihan $x_1, x_2,
\ldots, x_n$.  Jadi, jika $X_1,~X_2, \ldots,~X_n$ saling bebas, maka fungsi distribusi
kumulatif gabungan peubah acak ${\bar X} = (X_1, X_2, \ldots,
X_n)$ hanyalah hasil kali fungsi-fungsi distribusi kumulatif masing-masing.  Ketika dua peubah
acak saling bebas, secara lebih singkat kita mengatakan keduanya \emx {
bebas.}\index{kebebasan peubah\\ acak}
\end{definition}

Dengan menggunakan Persamaan~\ref{eq 4.4}, teorema berikut dapat dengan mudah ditunjukkan berlaku untuk peubah acak
kontinu yang saling bebas.

\begin{theorem}\label{thm 4.2}
Misalkan $X_1$,~$X_2$, \ldots,~$X_n$ adalah peubah acak kontinu dengan fungsi densitas
$f_1(x),~f_2(x), \ldots,~f_n(x)$.  Maka peubah-peubah acak ini \emx {saling
bebas} jika dan hanya jika
$$f(x_1, x_2, \ldots, x_n) = f_1(x_1)f_2(x_2) \cdots  f_n(x_n)$$    untuk sebarang pilihan $x_1,
x_2, \ldots, x_n$.  
\end{theorem}

Mari kita tinjau beberapa contoh.

\begin{example}
Dalam contoh ini, kita mendefinisikan tiga peubah acak, $X_1,\ X_2$, dan $X_3$.  Kita akan menunjukkan bahwa
$X_1$ dan $X_2$ saling bebas, sedangkan $X_1$ dan $X_3$ tidak saling bebas.  Pilihlah sebuah titik
$\omega = (\omega_1,\omega_2)$ secara acak dari persegi satuan.  Tetapkan $X_1 =
\omega_1^2$, $X_2 = \omega_2^2$, dan $X_3 =
\omega_1 + \omega_2$.  Carilah distribusi gabungan $F_{12}(r_1,r_2)$ dan $F_{23}(r_2,r_3)$.

Kita telah melihat (lihat Contoh~\ref{exam 2.2.7.1}) bahwa
\begin{eqnarray*}
 F_1(r_1) & = & P(-\infty < X_1 \leq r_1) \\
         & = & \sqrt{r_1}, \qquad \mbox{jika} \,\,0 \leq r_1 \leq 1\ ,
\end{eqnarray*}
dan dengan cara serupa,
$$ F_2(r_2) = \sqrt{r_2}\ ,$$
jika $0 \leq r_2 \leq 1$.
Sekarang kita mempunyai (lihat Gambar~\ref{fig 5.15})
\putfig{3truein}{PSfig5-15}{$X_1$ dan $X_2$ saling bebas.}{fig 5.15}
\begin{eqnarray*}
 F_{12}(r_1,r_2) & = & P(X_1 \leq r_1 \,\, \mbox{dan}\,\, X_2 \leq r_2) \\
                & = & P(\omega_1 \leq \sqrt{r_1} \,\,\mbox{dan}\,\, \omega_2 \leq
\sqrt{r_2}) \\
                & = & \mbox{Luas}\,(E_1)\\
                & = & \sqrt{r_1} \sqrt{r_2} \\
                & = &F_1(r_1)F_2(r_2)\ .
\end{eqnarray*}
Dalam kasus ini $F_{12}(r_1,r_2) = F_1(r_1)F_2(r_2)$ sehingga $X_1$ dan $X_2$
saling bebas.  Di sisi lain, jika $r_1 = 1/4$ dan $r_3 = 1$, maka (lihat Gambar~\ref{fig
5.16})
\putfig{3truein}{PSfig5-16}{$X_1$ dan $X_3$ tidak saling bebas.}{fig 5.16}
\begin{eqnarray*}
F_{13}(1/4,1) & = & P(X_1 \leq 1/4,\ X_3 \leq 1) \\
              & = & P(\omega_1 \leq 1/2,\ \omega_1 + \omega_2 \leq 1) \\
              & = & \mbox{Luas}\,(E_2) \\
              & = & \frac 12 - \frac 18 = \frac 38\ .
\end{eqnarray*}
Sekarang, dengan mengingat bahwa
$$F_3(r_3) = \left \{ \matrix{
                  0, & \mbox{jika} \,\,r_3 < 0, \cr
         (1/2)r_3^2, & \mbox{jika} \,\,0 \leq r_3 \leq 1, \cr
   1-(1/2)(2-r_3)^2, & \mbox{jika} \,\,1 \leq r_3 \leq 2, \cr
                  1, & \mbox{jika} \,\,2 < r_3,\cr}\right.
$$ (lihat Contoh~\ref{exam 2.2.7.2}), kita mempunyai $F_1(1/4)F_3(1) = (1/2)(1/2) = 1/4$.  Oleh karena itu, $X_1$
dan $X_3$ bukan peubah acak yang saling bebas.  Perhitungan serupa menunjukkan bahwa $X_2$ dan
$X_3$ juga tidak saling bebas.
\end{example}

Meskipun tidak akan kita buktikan di sini, teorema berikut berguna.  Pernyataan ini
juga berlaku bagi peubah acak diskret yang saling bebas.  Suatu bukti dapat
ditemukan dalam R\'enyi.\index{R\'ENYI, A.}\footnote{A. R\'enyi, \emx {Probability Theory} (Budapest:
Akad\'emiai Kiad\'o, 1970), p.~183.}

\begin{theorem}\label{thm 4.3}
Misalkan $X_1, X_2, \ldots, X_n$ adalah peubah acak kontinu yang saling bebas dan misalkan
$\phi_1(x), \phi_2(x), \ldots, \phi_n(x)$ adalah fungsi kontinu.  Maka $\phi_1(X_1),$
\newline
$\phi_2(X_2), \ldots, \phi_n(X_n)$ saling bebas.
\end{theorem}

\subsection*{Percobaan Bebas}

Dengan menggunakan gagasan kebebasan, sekarang kita dapat merumuskan gagasan
percobaan bebas untuk ruang sampel kontinu (lihat Definisi~\ref{def 5.5}).

\begin{definition}\label{def 5.12} Suatu barisan $X_1$,~$X_2$, \dots,~$X_n$ dari peubah acak
$X_i$ yang saling bebas dan memiliki densitas yang sama disebut \emx {proses
percobaan\linebreak[4] bebas.}\index{proses percobaan bebas}
\end{definition}
\par
Seperti pada peubah acak diskret, proses percobaan bebas ini muncul secara
alami dalam situasi ketika suatu percobaan yang dinyatakan oleh satu peubah acak diulang
$n$ kali.

\subsection*{Densitas Beta}

Selanjutnya kita meninjau contoh yang melibatkan ruang sampel dengan koordinat diskret
maupun kontinu.  Untuk contoh ini kita memerlukan fungsi densitas baru
yang disebut \emx {densitas beta.}\index{fungsi densitas!beta}\index{densitas beta}  Densitas ini
memiliki dua parameter,
$\alpha$,~$\beta$, dan didefinisikan oleh
$$
B(\alpha,\beta,x) =  \left \{ \matrix{ 
(1/B(\alpha,\beta))x^{\alpha - 1}(1 - x)^{\beta - 1}, & {\mbox{jika}}\,\, 0 \leq x \leq 1, \cr
                                                   0, & {\mbox{lainnya}}.\cr}\right. 
$$
Di sini $\alpha$ dan $\beta$ adalah sebarang bilangan positif, dan fungsi beta
$B(\alpha,\beta)$ diberikan oleh luas di bawah grafik $x^{\alpha - 1}(1 -
x)^{\beta - 1}$ antara 0 dan 1:
$$
B(\alpha,\beta) = \int_0^1 x^{\alpha - 1}(1 - x)^{\beta - 1}\,dx\ .
$$
Perhatikan bahwa ketika $\alpha = \beta = 1$, densitas beta adalah densitas seragam. 
Ketika $\alpha$ dan $\beta$ lebih besar dari 1, densitas tersebut berbentuk lonceng, tetapi
ketika keduanya kurang dari 1, densitas itu berbentuk U, sebagaimana ditunjukkan oleh contoh-contoh dalam
Gambar~\ref{fig 4.6}.

\putfig{4truein}{PSfig4-6}{Densitas beta untuk $\alpha = \beta = .5,1,2.$}{fig 4.6}

Kita hanya akan memerlukan nilai fungsi beta untuk nilai bulat
$\alpha$ dan $\beta$, dan dalam kasus ini
$$
B(\alpha,\beta) = \frac{(\alpha - 1)!\,(\beta - 1)!}{(\alpha + \beta - 1)!}\ .
$$

\begin{example}\label{exam 4.16}
Dalam masalah medis, sering diasumsikan bahwa suatu obat efektif dengan
peluang~$x$ setiap kali digunakan dan berbagai percobaan bersifat
saling bebas, sehingga pada hakikatnya kita melempar koin bias dengan peluang
$x$ untuk kepala.  
Sebelum melakukan percobaan lebih lanjut, Anda tidak mengetahui nilai $x$, tetapi pengalaman sebelumnya mungkin memberikan
sejumlah informasi tentang nilai-nilai yang mungkin. Wajar untuk menyatakan informasi ini
dengan membuat sketsa fungsi densitas
guna menentukan distribusi bagi $x$.  Jadi, kita memandang $x$ sebagai peubah acak
kontinu yang mengambil nilai antara 0 dan 1.  Jika Anda sama sekali tidak memiliki pengetahuan, Anda akan
membuat sketsa densitas seragam.  Jika pengalaman sebelumnya menunjukkan bahwa
$x$ sangat mungkin berada dekat 2/3, Anda akan membuat sketsa densitas dengan maksimum di 2/3 dan
sebaran yang mencerminkan ketidakpastian Anda dalam taksiran 2/3 tersebut. Kemudian Anda ingin mencari
fungsi densitas yang secara wajar sesuai dengan sketsa Anda.  Densitas beta menyediakan suatu kelas
densitas yang dapat disesuaikan dengan sebagian besar sketsa yang mungkin Anda buat. Sebagai contoh, untuk
$\alpha > 1$ dan $\beta > 1$, densitas itu berbentuk lonceng, dengan parameter $\alpha$ dan $\beta$
menentukan puncak dan sebarannya. 
\par
Anggaplah pelaku percobaan telah memilih suatu densitas beta untuk menyatakan keadaan
pengetahuannya tentang $x$ sebelum percobaan.  Kemudian ia memberikan obat tersebut kepada $n$
subjek dan mencatat banyaknya sukses, $i$.   Bilangan $i$ merupakan peubah acak
diskret, sehingga kita dapat dengan mudah menyatakan himpunan hasil yang mungkin dari percobaan ini
dengan merujuk pada pasangan terurut $(x, i)$.
\par
Misalkan $m(i|x)$ menyatakan peluang bahwa kita mengamati $i$ sukses dengan syarat nilai $x$.  Berdasarkan
asumsi kita, $m(i|x)$ merupakan distribusi binomial dengan peluang~$x$ untuk sukses:
$$
m(i|x) = b(n,x,i) = {n \choose i} x^i(1 - x)^j\ ,
$$
dengan $j = n - i$.  
\par
Jika $x$ dipilih secara acak dari $[0,1]$ dengan densitas beta
$B(\alpha,\beta,x)$, maka fungsi densitas bagi hasil berupa
pasangan $(x,i)$ adalah
\begin{eqnarray*}
f(x,i) & = & m(i|x)B(\alpha,\beta,x) \\
       & = & {n \choose i} x^i(1 - x)^j \frac 1{B(\alpha,\beta)} x^{\alpha - 1}(1 -
x)^{\beta - 1} \\
       & = & {n \choose i}  \frac 1{B(\alpha,\beta)} x^{\alpha + i - 1}(1 - x)^{\beta +
j - 1}\ .
\end{eqnarray*}
Sekarang misalkan $m(i)$ adalah peluang bahwa kita mengamati $i$ sukses \emx {tanpa}
mengetahui nilai $x$.  Maka
\begin{eqnarray*}
m(i) & = & \int_0^1 m(i|x) B(\alpha,\beta,x)\,dx \\
     & = & {n \choose i} \frac 1{B(\alpha,\beta)} \int_0^1 x^{\alpha + i - 1}(1 -
x)^{\beta + j - 1}\,dx \\
     & = & {n \choose i} \frac {B(\alpha + i,\beta + j)}{B(\alpha,\beta)}\ .
\end{eqnarray*}
Oleh karena itu, densitas peluang $f(x|i)$ untuk $x$, dengan syarat $i$ sukses
teramati, adalah
$$
f(x|i)  =  \frac {f(x,i)}{m(i)} 
$$
\begin{equation}
\ \ \ \ \ \ \ \ \ \ \ \ \ \ \ \ \ \ \ \ \ \ \ \ \ \ \ \ \ =  
\frac {x^{\alpha + i - 1}(1 - x)^{\beta + j - 1}}{B(\alpha + i,\beta
+ j)}\ ,\label{eq 4.5}
\end{equation}
yakni, $f(x|i)$ merupakan densitas beta lainnya.  Hal ini mengatakan bahwa jika kita mengamati $i$
sukses dan $j$ kegagalan pada $n$ subjek, maka densitas baru bagi
peluang bahwa obat tersebut efektif kembali berupa densitas beta, tetapi dengan
parameter $\alpha + i$, $\beta + j$.
\par
Sekarang kita menganggap bahwa sebelum percobaan, kita memilih densitas beta dengan
parameter $\alpha$ dan $\beta$, dan bahwa dalam percobaan tersebut
kita memperoleh $i$ sukses dalam $n$ percobaan.  Baru saja kita melihat bahwa dalam kasus ini,
densitas baru bagi $x$
adalah densitas beta dengan parameter $\alpha + i$ dan $\beta + j$.  
\par
Sekarang kita ingin menghitung peluang bahwa obat tersebut efektif pada subjek berikutnya.  Untuk setiap
bilangan real tertentu $t$ antara 0 dan 1, peluang bahwa $x$ memiliki nilai $t$ diberikan oleh
ekspresi dalam Persamaan~\ref{eq 4.5}.  Dengan syarat $x$ bernilai $t$, peluang bahwa obat
efektif pada subjek berikutnya hanyalah $t$.  Jadi, untuk memperoleh peluang bahwa obat tersebut efektif
pada subjek berikutnya, kita mengintegralkan hasil kali ekspresi dalam Persamaan~\ref{eq 4.5} dan $t$ atas semua
nilai $t$ yang mungkin.  Kita memperoleh:
\begin{eqnarray*}
\lefteqn{\frac 1{B(\alpha + i,\beta + j)} 
\int_0^1 t\cdot t^{\alpha + i - 1}(1 - t)^{\beta + j - 1}\,dt}
\\
  & = & \frac {B(\alpha + i + 1,\beta + j)}{B(\alpha + i,\beta + j)} \\
  & = & \frac {(\alpha + i)!\,(\beta + j - 1)!}{(\alpha + \beta + i + j)!} \cdot 
\frac {(\alpha + \beta + i + j - 1)!}{(\alpha + i - 1)!\,(\beta + j - 1)!} \\
  & = & \frac {\alpha + i}{\alpha + \beta + n}\ .
\end{eqnarray*}
Jika $n$ besar, maka taksiran kita bagi peluang sukses setelah percobaan
kira-kira sama dengan proporsi sukses yang diamati dalam percobaan, yang
tentu merupakan kesimpulan yang wajar.
\end{example}

Contoh berikut merupakan contoh lain ketika peluang yang sebenarnya tidak diketahui dan harus
ditaksir berdasarkan data percobaan.

\begin{example}(Masalah bandit berlengan dua)\label{exam 4.17}\index{bandit berlengan dua}
Anda berada di sebuah kasino dan berhadapan dengan dua mesin slot.  Setiap mesin membayar
1 dolar atau tidak membayar apa pun.  Peluang bahwa mesin pertama membayar
satu dolar adalah $x$, sedangkan peluang bahwa mesin kedua membayar satu dolar adalah $y$.  Kita menganggap
bahwa $x$ dan $y$ adalah bilangan acak yang dipilih secara saling bebas dari interval
$[0,1]$ dan tidak Anda ketahui.  Anda diizinkan bermain sepuluh kali,
dengan memilih salah satu mesin setiap kali.  Bagaimanakah Anda harus memilih agar
banyaknya kemenangan menjadi maksimum?

Salah satu strategi yang terdengar masuk akal adalah menghitung pada setiap tahap
peluang bahwa masing-masing mesin akan membayar, lalu memilih mesin dengan
peluang yang lebih tinggi.  Misalkan menang($i$), untuk~$i = 1$ atau~2, adalah berapa kali
Anda menang pada mesin ke-$i$.  Demikian pula, misalkan kalah($i$) adalah berapa kali
Anda kalah pada mesin ke-$i$.  Maka, dari Contoh~\ref{exam
4.16}, peluang $p(i)$ bahwa Anda menang jika memilih mesin ke-$i$ adalah 
$$
p(i) = \frac {{\mbox{menang}}(i) + 1} {{\mbox{menang}}(i) + {\mbox{kalah}}(i) + 2}\ .
$$
Jadi, jika $p(1) > p(2)$, Anda akan memainkan mesin 1; jika tidak, Anda akan memainkan
mesin 2.  Kita telah menulis program {\bf TwoArm}\index{TwoArm (program)} untuk menyimulasikan
percobaan ini.  Dalam program tersebut, pengguna menentukan nilai awal $x$ dan
$y$ (tetapi nilai-nilai ini tidak diketahui oleh pelaku percobaan).  Pada
setiap tahap, program menghitung dua densitas bersyarat bagi $x$ dan $y$ dengan syarat hasil
percobaan sebelumnya, lalu menghitung $p(i)$ untuk~$i = 1$,~2.  Kemudian program memilih
mesin dengan nilai peluang menang tertinggi untuk permainan
berikutnya.  Program mencetak mesin yang dipilih pada setiap permainan beserta hasil
permainan tersebut.  Program juga menggambar densitas baru bagi~$x$ (garis utuh) dan $y$
(garis putus-putus), dengan hanya menampilkan densitas saat ini.  Kita telah menjalankan program selama
sepuluh permainan untuk kasus $x = .6$ dan $y = .7$.  Hasilnya ditampilkan dalam
Gambar~\ref{fig 4.7}.

\putfig{4.5truein}{PSfig4-7}{Mainkan mesin terbaik.}{fig 4.7}

Jalannya program menunjukkan kelemahan strategi ini.  Peluang awal kita untuk
menang pada mesin yang lebih baik di antara keduanya adalah .7.  Kita mulai dengan
mesin yang lebih buruk, dan hasil-hasilnya membuat kita selalu memiliki peluang menang
lebih besar dari .6, sehingga kita terus memainkan mesin ini walaupun
mesin lainnya lebih baik.  Seandainya kalah pada permainan pertama, kita akan
berpindah mesin.  Densitas akhir kita bagi $y$ sama dengan densitas
awal, yakni densitas seragam.  Densitas akhir kita bagi $x$ berbeda
dan mencerminkan pengetahuan yang jauh lebih akurat tentang $x$.  Komputer bekerja cukup
baik dengan strategi ini dan menang pada tujuh dari sepuluh percobaan, tetapi sepuluh percobaan
belum cukup untuk menilai apakah strategi ini baik dalam jangka panjang.

\putfig{4.5truein}{PSfig4-8}{Mainkan pemenang.}{fig 4.8}

Strategi populer lainnya adalah \emx {strategi mainkan-pemenang.}  Seperti tersirat dari namanya,
dalam strategi ini kita memilih mesin yang sama ketika menang dan berpindah
mesin ketika kalah.  Program {\bf TwoArm} juga akan menyimulasikan
strategi ini.  Dalam Gambar~\ref{fig 4.8}, kita memperlihatkan hasil menjalankan program
dengan strategi mainkan-pemenang dan peluang sebenarnya yang sama, yaitu .6 dan .7,
untuk kedua mesin.  Setelah sepuluh permainan, densitas kita bagi peluang
menang yang tidak diketahui menunjukkan bahwa mesin kedua memang
lebih baik di antara keduanya.  Kita kembali menang pada tujuh dari sepuluh percobaan.

Tidak satu pun strategi yang kita simulasikan merupakan strategi terbaik untuk
memaksimumkan rata-rata kemenangan kita.  Strategi terbaik tersebut sangat rumit, tetapi
dihampiri secara wajar oleh strategi mainkan-pemenang.  Berbagai variasi
contoh ini berperan penting dalam masalah uji klinis obat,
ketika pelaku percobaan menghadapi situasi serupa.
\end{example}

\exercises
\begin{LJSItem}

\i\label{exer 4.2.1} Pilih sebuah titik $x$ secara acak (dengan densitas seragam) di dalam interval
$[0,1]$.  Tentukan peluang bahwa $x > 1/2$, dengan syarat
\begin{enumerate}
\item $x > 1/4$.

\item $x < 3/4$.

\item $|x - 1/2| < 1/4$.

\item $x^2 - x + 2/9 < 0$.
\end{enumerate}

\i\label{exer 4.2.2} Suatu bahan radioaktif memancarkan partikel $\alpha$ dengan laju yang dijelaskan oleh
fungsi densitas
$$
f(t) = .1e^{-.1t}\ .
$$
Tentukan peluang
bahwa sebuah partikel dipancarkan dalam 10~detik pertama, dengan syarat
\begin{enumerate}
\item tidak ada partikel yang dipancarkan dalam detik pertama.

\item tidak ada partikel yang dipancarkan dalam 5 detik pertama.

\item sebuah partikel dipancarkan dalam 3 detik pertama.

\item sebuah partikel dipancarkan dalam 20 detik pertama.
\end{enumerate}

\i\label{exer 4.2.3} Bola lampu Acme Super\index{bola lampu} diketahui memiliki masa guna yang
dijelaskan oleh fungsi densitas
$$
f(t) = .01e^{-.01t}\ ,
$$
dengan waktu $t$ diukur dalam jam.  
\begin{enumerate}
\item Tentukan \emx {laju kegagalan} bola lampu ini 
(lihat Latihan~\ref{sec 2.2}.\ref{exer 2.2.6}).

\item Tentukan \emx {keandalan} bola lampu ini setelah 20 jam.

\item Dengan syarat bola lampu bertahan selama 20 jam, tentukan peluang bahwa bola lampu
tersebut bertahan selama 20 jam lagi.

\item Tentukan peluang bahwa bola lampu padam pada jam keempat puluh satu, dengan syarat
bola lampu bertahan selama 40 jam.
\end{enumerate}

\i\label{exer 4.2.4} Misalkan Anda melempar anak panah pada sasaran berbentuk lingkaran berjari-jari 10 inci.  Dengan syarat
anak panah jatuh di setengah bagian atas sasaran, tentukan peluang bahwa
\begin{enumerate}
\item anak panah jatuh di setengah bagian kanan sasaran.

\item jaraknya dari pusat kurang dari 5 inci.

\item jaraknya dari pusat lebih dari 5 inci.

\item anak panah jatuh dalam jarak 5 inci dari titik $(0,5)$.
\end{enumerate}

\i\label{exer 4.2.5} Misalkan Anda memilih dua bilangan $x$ dan $y$, secara acak dan saling bebas dari
interval $[0,1]$.  Dengan syarat jumlah keduanya terletak dalam interval $[0,1]$,
tentukan peluang bahwa
\begin{enumerate}
\item $|x - y| < 1$.

\item $xy < 1/2$.

\item $\max\{x,y\} < 1/2$.

\item $x^2 + y^2 < 1/4$.

\item $x > y$.
\end{enumerate}

\i\label{exer 4.2.6} Tentukan fungsi-fungsi densitas bersyarat untuk percobaan berikut.
\begin{enumerate}
\item Sebuah bilangan $x$ dipilih secara acak dalam interval $[0,1]$, dengan syarat
$x > 1/4$.

\item Sebuah bilangan $t$ dipilih secara acak dalam interval $[0,\infty)$ dengan
densitas eksponensial $e^{-t}$, dengan syarat $1 < t < 10$.

\item Sebuah anak panah dilemparkan pada sasaran berbentuk lingkaran berjari-jari 10 inci, dengan syarat
anak panah itu jatuh di setengah bagian atas sasaran.

\item Dua bilangan $x$ dan $y$ dipilih secara acak dalam interval
$[0,1]$, dengan syarat $x > y$.
\end{enumerate}

\i\label{exer 4.2.7} Misalkan $x$ dan $y$ dipilih secara acak dari interval $[0,1]$.  Tunjukkan
bahwa kejadian $x > 1/3$ dan $y > 2/3$ merupakan kejadian yang saling bebas.

\i\label{exer 4.2.8} Misalkan $x$ dan $y$ dipilih secara acak dari interval $[0,1]$.  Pasangan
manakah di antara kejadian berikut yang saling bebas?
\begin{enumerate}
\item $x > 1/3$.

\item $y > 2/3$.

\item $x > y$.

\item $x + y < 1$.
\end{enumerate}

\i\label{exer 4.2.8.5}  Misalkan $X$ dan $Y$ adalah peubah acak kontinu dengan
fungsi densitas masing-masing $f_X(x)$ dan $f_Y(y)$.  Misalkan $f(x, y)$ menyatakan fungsi
densitas gabungan dari $(X, Y)$.  Tunjukkan bahwa
$$
\int_{-\infty}^\infty f(x, y)\, dy = f_X(x)\ ,
$$ dan
$$
\int_{-\infty}^\infty f(x, y)\, dx = f_Y(y)\ .
$$


\istar\label{exer 4.2.9} Dalam Latihan~\ref{sec 2.2}.\ref{exer 2.2.13} Anda
telah membuktikan hal berikut: Jika Anda mengambil sebatang tongkat sepanjang satu satuan dan mematahkannya menjadi tiga bagian,
dengan memilih titik-titik patahnya secara acak (yakni, memilih dua bilangan real secara saling bebas 
dan seragam dari [0, 1]), maka peluang bahwa ketiga bagian tersebut membentuk
sebuah segitiga adalah 1/4.  Sekarang tinjau percobaan serupa: Pertama patahkan tongkat secara
acak, lalu patahkan bagian yang lebih panjang secara acak.  Tunjukkan bahwa kedua percobaan
tersebut sebenarnya cukup berbeda, sebagai berikut:
\begin{enumerate}
\item Tulis sebuah program yang menyimulasikan kedua kasus dalam satu rangkaian 1000 percobaan,
mencetak proporsi sukses untuk setiap rangkaian, dan mengulangi proses ini
sepuluh kali.  (Sebut sebuah percobaan sukses jika ketiga bagian tersebut memang membentuk segitiga.) 
Mintalah program Anda memilih $(x,y)$ secara acak dalam persegi satuan, dan dalam setiap kasus
gunakan $x$ dan $y$ untuk menentukan kedua titik patah.  Untuk setiap percobaan, mintalah program itu menggambar
$(x,y)$ jika $(x,y)$ menghasilkan sukses.

\item Tunjukkan bahwa dalam percobaan kedua, peluang sukses teoretis
sebenarnya adalah $2\log 2 - 1$.
\end{enumerate}

\i\label{exer 4.2.10} Sebuah koin memiliki bias tak diketahui $p$ yang diasumsikan
berdistribusi seragam antara 0 dan 1.  Koin tersebut dilempar $n$ kali dan sisi kepala muncul
$j$ kali serta sisi ekor muncul $k$ kali.  Kita telah melihat bahwa peluang
sisi kepala muncul pada lemparan berikutnya adalah
$$
\frac {j + 1}{n + 2}\ .
$$
Tunjukkan bahwa nilai ini sama dengan peluang bahwa bola berikutnya berwarna hitam dalam
model guci Polya\index{model guci Polya} pada Latihan~\ref{sec 4.1}.\ref{exer 4.1.18}. 
Gunakan hasil ini untuk menjelaskan mengapa, dalam model guci Polya, proporsi bola hitam
tidak cenderung menuju 0~atau~1 seperti yang mungkin diduga, tetapi justru menuju
distribusi seragam pada interval $[0,1]$.

\i\label{exer 4.2.11} Pengalaman sebelumnya dengan suatu obat menunjukkan bahwa peluang~$p$
obat tersebut efektif merupakan besaran acak yang memiliki densitas beta dengan
parameter $\alpha = 2$ dan $\beta = 3$.  Obat itu digunakan pada sepuluh subjek dan
terbukti berhasil pada empat dari sepuluh subjek tersebut.  Densitas apa yang sekarang
harus kita tetapkan bagi peluang~$p$?  Berapa peluang bahwa obat itu
akan berhasil pada penggunaan berikutnya?

\i\label{exer 4.2.12} Tulis sebuah program yang memungkinkan Anda membandingkan strategi mainkan-pemenang
dan mainkan-mesin-terbaik untuk masalah bandit berlengan dua pada Contoh~\ref{exam
4.17}.  Mintalah program Anda menentukan peluang awal bahwa setiap mesin membayar
dengan memilih sepasang bilangan acak antara 0 dan 1.  Mintalah
program Anda menjalankan 20 permainan dan mencatat banyaknya kemenangan untuk masing-masing dari
kedua strategi.  Terakhir, mintalah program Anda melakukan 1000 pengulangan rangkaian 20
permainan dan menghitung rata-rata kemenangan per 20 permainan.  Strategi mana yang tampaknya
terbaik?  Ulangi simulasi ini dengan mengganti 20 menjadi 100.  Apakah jawaban Anda atas
pertanyaan di atas berubah?

\i\label{exer 4.2.13} Tinjau masalah bandit berlengan dua pada Contoh~\ref{exam 4.17}. 
Bruce Barnes\index{BARNES, B.} mengusulkan strategi berikut, yang merupakan variasi dari
strategi mainkan-mesin-terbaik.  Mesin dengan peluang
menang terbesar dimainkan \emx {kecuali jika} dua syarat berikut terpenuhi: (a)~selisih
peluang menang kurang dari .08, dan (b)~rasio
banyaknya permainan pada mesin yang lebih sering dimainkan terhadap banyaknya
permainan pada mesin yang lebih jarang dimainkan lebih besar dari 1.4.  Jika
kedua syarat di atas terpenuhi, mesin dengan peluang
menang yang lebih kecil dimainkan.  Tulis sebuah program untuk menyimulasikan strategi ini.  Mintalah
program Anda memilih peluang pembayaran awal setiap mesin secara acak dari
interval satuan $[0,1]$, menjalankan 20 permainan, dan mencatat banyaknya kemenangan. 
Ulangi percobaan ini 1000 kali dan peroleh rata-rata banyaknya kemenangan per 20
permainan.  Terapkan strategi kedua---misalnya, strategi mainkan-mesin-terbaik atau strategi
pilihan Anda sendiri, lalu lihat bagaimana strategi kedua ini dibandingkan dengan strategi Bruce berdasarkan
rata-rata banyaknya kemenangan.

%\end{LJSItem}}
\end{LJSItem}

\section{Paradoks}\label{sec 4.3}

Sebagian besar bagian ini didasarkan pada artikel Snell dan Vanderbei.\index{SNELL, J.
L.}\index{VANDERBEI, R.}\footnote{J. L. Snell and R. Vanderbei, ``Three Bewitching
Paradoxes," in \emx {Topics in Contemporary Probability and Its Applications}, CRC Press, Boca
Raton, 1995.}
\par
Kita harus sangat berhati-hati ketika menangani persoalan yang melibatkan peluang
bersyarat.  Pembaca tentu ingat bahwa dalam masalah Monty Hall
(Contoh~\ref{exam 4.5}), jika peserta memilih pintu yang menyembunyikan mobil, Monty
dapat memilih pintu mana yang akan dibuka.  Kita membuat asumsi bahwa dalam hal ini,
ia memilih setiap pintu dengan peluang 1/2.  Kita kemudian mencatat bahwa jika asumsi
ini diubah, jawaban atas pertanyaan semula juga berubah.  Dalam bagian ini, kita akan
mempelajari contoh-contoh lain dari fenomena yang sama.

\begin{example}\label{exam 4.3.1}
Tinjau sebuah keluarga dengan dua anak.  Dengan diketahui bahwa salah satu anak
adalah laki-laki, berapakah peluang bahwa kedua anak adalah laki-laki?
\par
Salah satu cara mendekati masalah ini adalah mengatakan bahwa anak lainnya mempunyai
peluang sama untuk menjadi laki-laki atau perempuan, sehingga peluang kedua anak
laki-laki adalah 1/2.  Penyelesaian ``buku teks" akan menggambar diagram pohon, lalu
membentuk pohon bersyarat dengan menghapus lintasan sehingga hanya lintasan yang
sesuai dengan informasi yang diberikan yang tersisa.  Hasilnya ditampilkan dalam
Gambar~\ref{fig 4.3}. Kita melihat bahwa peluang dua anak laki-laki dengan diketahui
ada anak laki-laki dalam keluarga bukanlah 1/2, melainkan 1/3.
\putfig{3.5truein}{PSfig4-3}{Pohon untuk Contoh \protect\ref{exam 4.3.1}\protect.}{fig 4.3}
\end{example}

Masalah ini dan masalah sejenis dibahas oleh Bar-Hillel dan
Falk.\index{BAR-HILLEL, M.}\index{FALK, R.}\footnote{M. Bar-Hillel and R. Falk, ``Some teasers
concerning conditional probabilities," \emx {Cognition}, vol. 11 (1982), pgs. 109-122.} Para
penulis menekankan bahwa jawaban atas peluang bersyarat semacam ini dapat berubah,
bergantung pada cara informasi yang diberikan benar-benar diperoleh.  Sebagai contoh,
mereka menunjukkan bahwa 1/2 merupakan jawaban yang benar untuk skenario berikut.

\begin{example}\label{exam 4.3.2}
Tn.\ Smith memiliki dua anak.  Kita bertemu dengannya ketika ia berjalan di jalan
bersama seorang anak laki-laki yang dengan bangga ia perkenalkan sebagai putranya.
Berapakah peluang bahwa anak Tn.\ Smith yang lain juga laki-laki?
\par
Seperti biasa, kita harus membuat beberapa asumsi tambahan.  Sebagai contoh, kita
akan mengasumsikan bahwa jika Tn.\ Smith memiliki seorang anak laki-laki dan seorang
anak perempuan, peluang ia memilih salah satunya untuk menemaninya berjalan adalah
sama. Dalam Gambar~\ref{fig 4.3.5}, kita menampilkan analisis pohon bagi masalah ini
dan melihat bahwa 1/2 memang merupakan jawaban yang benar.
\putfig{5truein}{PSfig4-3-5}{Pohon untuk Contoh \protect\ref{exam 4.3.2}\protect.}{fig 4.3.5}
\end{example}

\begin{example}\label{exam 4.3.3}
Tidak mudah memikirkan skenario masuk akal yang menghasilkan jawaban klasik 1/3.
Stephen Geller\index{GELLER, S.} mencoba melakukannya ketika mengajukan masalah ini kepada Marilyn vos
Savant.\index{vos SAVANT, M.}\footnote{M. vos Savant, ``Ask Marilyn," \emx {Parade
Magazine}, 9  September;  2 December;  17  February 1990, reprinted in Marilyn
vos Savant, \emx {Ask Marilyn}, St. Martins, New York, 1992.}
Masalah Geller adalah sebagai berikut: Seorang pemilik toko berkata bahwa ia memiliki
dua anak anjing beagle baru untuk diperlihatkan kepada Anda, tetapi ia tidak tahu
apakah keduanya jantan, keduanya betina, atau masing-masing satu.  Anda mengatakan
bahwa Anda hanya menginginkan yang jantan, lalu ia menelepon orang yang sedang
memandikan keduanya.  ``Apakah setidaknya satu jantan?" tanyanya. ``Ya," katanya
kepada Anda sambil tersenyum.  Berapakah peluang bahwa anak anjing yang \emx {lain}
juga jantan?
\par
Pembaca diminta memutuskan apakah model yang memberikan jawaban 1/3 layak digunakan
dalam kasus ini.
\end{example}

Dalam contoh-contoh sebelumnya, paradoks yang tampak dapat diselesaikan dengan mudah
dengan menyatakan secara jelas model yang digunakan dan asumsi yang dibuat.  Sekarang
kita beralih ke beberapa contoh yang paradoksnya tidak begitu mudah diselesaikan.

\begin{example}\label{exam 4.3.4}
Dua amplop masing-masing berisi sejumlah uang.  Satu amplop diberikan kepada
Ali\index{Ali} dan yang lain kepada Baba\index{Baba}; mereka diberi tahu bahwa satu
amplop berisi uang dua kali lebih banyak daripada amplop lainnya.  Namun, keduanya
tidak tahu siapa yang mendapat hadiah lebih besar.  Sebelum ada yang membuka amplop,
Ali ditanya apakah ia ingin menukar amplopnya dengan Baba.  Ali bernalar sebagai
berikut: Misalkan jumlah dalam amplop saya adalah $x$.  Jika saya bertukar, saya akan
memperoleh $x/2$ dengan peluang 1/2 dan $2x$ dengan peluang 1/2.  Jika saya diberi
kesempatan memainkan permainan ini berkali-kali dan selalu bertukar, secara rata-rata
saya akan memperoleh
$$
\frac 12 \frac x2 + \frac 12 2x = \frac 54 x\ .
$$
Nilai ini lebih besar daripada kemenangan rata-rata saya jika saya tidak bertukar.
\par
Tentu saja, Baba mendapat kesempatan yang sama dan bernalar dengan cara serupa hingga
menyimpulkan bahwa ia juga ingin bertukar.  Maka mereka bertukar, dan masing-masing
berpikir bahwa kekayaan bersihnya baru saja meningkat 25\%.
\par
Karena keduanya belum membuka amplop, proses ini dapat diulangi dan mereka pun
bertukar lagi.  Sekarang mereka kembali memegang amplop semula, tetapi mereka berpikir
bahwa kekayaan masing-masing telah meningkat 25\% sebanyak dua kali.  Dengan penalaran
ini, mereka dapat meyakinkan diri bahwa dengan bertukar amplop berulang kali, mereka
dapat menjadi kaya tanpa batas.  Jelas ada yang salah dalam penalaran di atas, tetapi
di manakah kesalahannya?
\par
Salah satu siasat dalam membuat paradoks adalah menjadikannya sedikit lebih rumit
daripada yang diperlukan agar kita makin bingung.  Seperti yang disarankan John
Finn\index{FINN, J.}, paradoks ini sebenarnya dapat dimulai dengan masalah yang lebih
sederhana.  Misalkan Ali dan Baba tahu bahwa saya akan memberikan sebuah amplop
berisi \$5 atau sebuah amplop berisi \$10.  Saya akan melempar koin untuk menentukan
amplop mana yang diberikan kepada Ali, lalu memberikan amplop lainnya kepada Baba.
Ali kemudian dapat berargumen bahwa Baba memiliki $2x$ dengan peluang $1/2$ dan $x/2$
dengan peluang $1/2$.  Hal ini membawa Ali pada kesimpulan yang sama seperti
sebelumnya.  Namun, sekarang jelas bahwa penalaran ini tidak masuk akal. Jika Ali
memiliki amplop berisi \$5, Baba tidak mungkin memiliki setengahnya, yakni \$2.50,
karena jumlah itu bahkan bukan salah satu pilihan.  Demikian pula, jika Ali memiliki
\$10, Baba tidak mungkin memiliki dua kali lipatnya, yakni \$20. Dalam masalah yang
lebih sederhana ini, hasil-hasil yang mungkin ditampilkan oleh diagram pohon dalam
Gambar~\ref{fig 4.9}.
\putfig{3truein}{PSfig4-9}{Versi John Finn dari Contoh \protect\ref{exam 4.3.4}\protect.}{fig 4.9}%%4.5truein 
Dari diagram tersebut, jelas bahwa pertukaran tidak membuat salah satu pihak lebih untung.
\end{example}

Dalam contoh di atas, penalaran Ali keliru karena ia menyimpulkan bahwa jika jumlah
dalam amplopnya adalah $x$, peluang amplopnya memuat jumlah yang lebih kecil adalah
1/2, dan peluang amplopnya memuat jumlah yang lebih besar juga 1/2.  Sebenarnya,
peluang bersyarat ini bergantung pada distribusi jumlah uang yang dimasukkan ke dalam
amplop.
\par
Untuk memperjelas, misalkan $X$ adalah peubah acak bernilai bilangan bulat positif
yang menyatakan jumlah lebih kecil di antara dua jumlah dalam amplop.  Selain itu,
misalkan distribusi $X$ diberikan; yakni, untuk setiap bilangan bulat positif $x$,
kita diberi nilai
$$p_x = P(X = x)\ .$$
(Dalam contoh Finn, $p_5 = 1$, dan $p_{n} = 0$ untuk semua nilai $n$ lainnya.)
Kita lalu dapat menghitung dengan mudah peluang bersyarat bahwa sebuah amplop memuat
jumlah lebih kecil, dengan diketahui bahwa amplop tersebut memuat $x$ dolar.  Dua
titik sampel yang mungkin adalah $(x, x/2)$ dan $(x, 2x)$.  Jika $x$ ganjil, titik
sampel pertama memiliki peluang 0 karena $x/2$ bukan bilangan bulat, sehingga peluang
bersyarat yang dicari adalah 1 bahwa $x$ merupakan jumlah lebih kecil.  Jika $x$
genap, kedua titik sampel masing-masing memiliki peluang $p_{x/2}$ dan $p_x$, sehingga
peluang bersyarat bahwa $x$ merupakan jumlah lebih kecil adalah
$$\frac{p_x}{p_{x/2} + p_x}\ ,$$
yang tidak harus sama dengan 1/2.
\par
Steven Brams\index{BRAMS, S.} dan D. Marc Kilgour\index{KILGOUR, D. M.}\footnote{S. J. Brams
and D. M. Kilgour, ``The Box Problem:  To Switch or Not to Switch," \emx {Mathematics Magazine},
vol. 68, no. 1 (1995),  p. 29.} mempelajari, untuk berbagai distribusi, persoalan
apakah seseorang sebaiknya menukar amplop jika tujuannya memaksimumkan kemenangan
rata-rata jangka panjang.  Misalkan $x$ adalah jumlah dalam amplop Anda.  Mereka
menunjukkan bahwa untuk setiap distribusi $X$, terdapat setidaknya satu nilai $x$ yang
membuat Anda sebaiknya bertukar.  Mereka memberikan contoh distribusi yang hanya
memiliki satu nilai $x$ yang membuat Anda sebaiknya bertukar (lihat
Latihan~\ref{exer 4.3.5}).  Barangkali kasus yang paling menarik adalah distribusi
yang membuat Anda selalu sebaiknya bertukar.  Sekarang kita berikan contoh tersebut.

\begin{example}\label{exam 4.3.5}
Misalkan ada dua amplop\index{amplop} di hadapan kita, dan satu amplop berisi uang dua
kali lebih banyak daripada yang lain (kedua jumlah berupa bilangan bulat positif).
Kita diberi salah satu amplop dan ditanya apakah kita ingin bertukar.
\par
Seperti di atas, misalkan $X$ menyatakan jumlah lebih kecil di antara kedua jumlah
dalam amplop, dan misalkan
$$p_x = P(X = x)\ .$$
Sekarang kita dapat menghitung kemenangan rata-rata jangka panjang jika kita
bertukar.  (Rata-rata jangka panjang ini merupakan contoh konsep peluang yang dikenal
sebagai nilai harapan dan akan dibahas dalam Bab~\ref{chp 6}.)  Dengan diketahui bahwa
salah satu dari dua titik sampel telah terjadi, peluang bahwa titik tersebut adalah
$(x, x/2)$ adalah
$$\frac{p_{x/2}}{p_{x/2} + p_x}\ ,$$
dan peluang bahwa titik tersebut adalah $(x, 2x)$ ialah
$$\frac{p_x}{p_{x/2} + p_x}\ .$$
Jadi, jika kita bertukar, kemenangan rata-rata jangka panjang kita adalah
$$\frac{p_{x/2}}{p_{x/2} + p_x}\frac x2 + \frac{p_x}{p_{x/2} + p_x} 2x\ .$$
Jika nilai ini lebih besar daripada $x$, dalam jangka panjang kita memperoleh
keuntungan dengan bertukar.  Perhitungan aljabar rutin menunjukkan bahwa ekspresi di
atas lebih besar daripada $x$ jika dan hanya jika
\begin{equation}
\frac{p_{x/2}}{p_{x/2} + p_x} < \frac 23\ .
\label{eq 4.3.1}
\end{equation}
\par
Menarik untuk meninjau apakah terdapat distribusi pada bilangan bulat positif yang
membuat pertidaksamaan~\ref{eq 4.3.1} benar untuk semua nilai genap $x$.  Brams dan
Kilgour\footnote{ibid.} memberikan contoh berikut.
\par
Kita definisikan $p_x$ sebagai berikut:
$$
p_x = \left \{ \matrix{
               \frac 13 \Bigl(\frac 23\Bigr)^{k-1}, & \mbox{jika}\,\, x = 2^k, \cr
               0, & \mbox{selainnya.}\cr
}\right.
$$
Mudah dihitung (lihat Latihan~\ref{exer 4.3.4}) bahwa untuk semua nilai $x$ yang
relevan, kita mempunyai
$$
\frac{p_{x/2}}{p_{x/2} + p_x} = \frac 35\ ,
$$
yang berarti pertidaksamaan~\ref{eq 4.3.1} selalu benar.
\end{example}

Sejauh ini, kita dapat menyelesaikan paradoks dengan menyatakan secara jelas asumsi
yang dibuat dan secara tepat model yang digunakan.  Kita mengakhiri bagian ini dengan
menjelaskan paradoks yang tidak dapat kita selesaikan.

\begin{example}\label{exam 4.3.6}
Misalkan ada dua amplop\index{amplop} di hadapan kita, dan kita diberi tahu bahwa
kedua amplop masing-masing berisi $X$ dan $Y$ dolar, dengan $X$ dan $Y$ bilangan bulat
positif yang berbeda.  Kita memilih salah satu amplop secara acak dan membukanya;
misalkan isinya $X$.  Apakah mungkin menentukan, dengan peluang lebih besar daripada
1/2, apakah $X$ merupakan jumlah dolar yang lebih kecil?
\par
Bahkan jika kita tidak mengetahui distribusi bersama $X$ dan $Y$, secara mengejutkan
jawabannya adalah ya!  Caranya sebagai berikut.  Lemparkan koin adil hingga kepala
muncul untuk pertama kalinya.  Misalkan $Z$ menyatakan banyaknya lemparan yang
diperlukan ditambah 1/2.  Jika $Z > X$, kita mengatakan bahwa $X$ merupakan jumlah
yang lebih kecil, dan jika $Z < X$, kita mengatakan bahwa $X$ merupakan jumlah yang
lebih besar.
\par
Pertama, jika $Z$ berada di antara $X$ dan $Y$, kita pasti benar.  Karena $X$ dan $Y$
tidak sama, $Z$ berada di antara keduanya dengan peluang positif.  Kedua, jika $Z$
tidak berada di antara $X$ dan $Y$, maka $Z$ lebih besar daripada $X$ dan $Y$, atau
lebih kecil daripada $X$ dan $Y$.  Dalam kedua kasus, berdasarkan simetri, $X$ merupakan
jumlah yang lebih kecil dengan peluang 1/2 (ingat, kita memilih amplop secara acak).
Jadi, peluang bahwa kita benar lebih besar daripada 1/2.
\end{example}

\exercises
\begin{LJSItem}

\i\label{exer 4.3.1}
Salah satu paradoks peluang bersyarat yang pertama dikemukakan oleh
Bertrand.\index{BERTRAND, J.}\footnote{J. Bertrand, \emx {Calcul des Probabilit\'{e}s},
Gauthier-Uillars, 1888.}  Paradoks ini disebut \emx {Paradoks Kotak}\index{Paradoks
kotak}.  Sebuah lemari memiliki tiga laci. Laci pertama berisi dua bola emas, laci
kedua berisi dua bola perak, dan laci ketiga berisi satu bola perak serta satu bola
emas. Sebuah laci dipilih secara acak, lalu satu bola dipilih secara acak dari dua
bola di dalam laci.  Dengan diketahui bahwa bola emas terambil, berapakah peluang
bahwa laci yang berisi dua bola emaslah yang dipilih?

\i\label{exer 4.3.2}
Masalah berikut disebut \emx {masalah dua as}\index{masalah dua as}. Masalah yang
berasal dari tahun 1936 ini dikaitkan dengan matematikawan Inggris J. H. C.
Whitehead\index{WHITEHEAD, J. H. C.} (lihat Gridgeman\index{GRIDGEMAN, N. T.}\footnote{N. T. 
Gridgeman, Letter, \emx {American Statistician}, 21 (1967), pgs. 38-39.}). Masalah ini juga
diajukan kepada Marilyn vos Savant\index{vos SAVANT, M.} oleh pakar teka-teki matematika Martin
Gardner\index{GARDNER, M.}, yang menyebutnya sebagai salah satu masalah kesukaannya.
\par
Satu tangan bridge\index{bridge} telah dibagikan, yakni tiga belas kartu dibagikan
kepada setiap pemain.  Dengan diketahui bahwa pasangan Anda memiliki setidaknya satu
as, berapakah peluang bahwa ia memiliki setidaknya dua as?  Dengan diketahui bahwa
pasangan Anda memiliki as hati, berapakah peluang bahwa ia memiliki setidaknya dua as?
Jawablah pertanyaan ini untuk versi bridge yang menggunakan delapan kartu, yaitu empat
as dan empat raja, dengan dua kartu dibagikan kepada setiap pemain.  (Pembaca mungkin
ingin menyelesaikan masalah ini dengan setumpuk 52 kartu.)

\i\label{exer 4.3.3}
Dalam latihan sebelumnya, wajar untuk bertanya, ``Bagaimana kita memperoleh informasi
bahwa tangan tersebut memiliki sebuah as?" Gridgeman meninjau dua cara berbeda untuk
memperoleh informasi ini.  (Sekali lagi, asumsikan tumpukan terdiri atas delapan kartu.)
\begin{enumerate}
\item
Asumsikan bahwa orang yang memegang tangan ditanya, ``Sebutkan satu as dalam tangan
Anda," dan menjawab, ``As hati."  Berapakah peluang bahwa ia memiliki as kedua?
\item
Misalkan orang yang memegang tangan diberi pertanyaan yang lebih langsung, ``Apakah
Anda memiliki as hati?" dan jawabannya ya.  Berapakah peluang bahwa ia memiliki as kedua?
\end{enumerate}

\i\label{exer 4.3.4}  Dengan menggunakan notasi yang diperkenalkan dalam
Contoh~\ref{exam 4.3.5}, tunjukkan bahwa dalam contoh Brams\index{BRAMS, S.} dan
Kilgour\index{KILGOUR, D. M.}, jika $x$ merupakan pangkat positif dari 2, maka
$$
\frac{p_{x/2}}{p_{x/2} + p_x} = \frac 35\ .
$$

\i\label{exer 4.3.5} Dengan menggunakan notasi yang diperkenalkan dalam
Contoh~\ref{exam 4.3.5}, misalkan
$$
p_x = \left \{ \matrix{
               \frac 23 \Bigl(\frac 13\Bigr)^k, & \mbox{jika}\,\, x = 2^k, \cr
               0, & \mbox{selainnya.}\cr
}\right.
$$
Tunjukkan bahwa terdapat tepat satu nilai $x$ sedemikian sehingga jika amplop Anda
berisi $x$, Anda sebaiknya bertukar.


\istar\label{exer 4.3.6} (Hanya untuk pemain bridge\index{bridge}.  Dari
Sutherland.\index{SUTHERLAND, E.}\footnote{E. Sutherland, ``Restricted Choice\index{restricted
choice, principle of} --- Fact or Fiction?", \emx {Canadian Master Point}, November 1, 1993.})  
Misalkan kita adalah deklaran dalam satu tangan bridge dan memiliki raja, 9, 8, 7,
serta 2 dari suatu jenis kartu, sedangkan dummy memiliki as, 10, 5, dan 4 dari jenis
yang sama.  Misalkan kita ingin memainkan jenis kartu ini sedemikian rupa sehingga
peluang tidak mengalami kekalahan pada jenis tersebut menjadi maksimum.  Kita mulai
dengan memainkan 2 menuju as dan melihat bahwa ratu jatuh di sebelah kiri kita.  Kita
kemudian memainkan 10 dari dummy, dan lawan di sebelah kanan memainkan enam (setelah
memainkan tiga pada putaran pertama).  Haruskah kita melakukan finesse atau bermain
dengan mengandalkan jatuhnya kartu?


\end{LJSItem}
